% ****** Start of file apssamp.tex ******
%
%   This file is part of the APS files in the REVTeX 4.2 distribution.
%   Version 4.2a of REVTeX, December 2014
%
%   Copyright (c) 2014 The American Physical Society.
%
%   See the REVTeX 4 README file for restrictions and more information.
%
% TeX'ing this file requires that you have AMS-LaTeX 2.0 installed
% as well as the rest of the prerequisites for REVTeX 4.2
%
% See the REVTeX 4 README file
% It also requires running BibTeX. The commands are as follows:
%
%  1)  latex apssamp.tex
%  2)  bibtex apssamp
%  3)  latex apssamp.tex
%  4)  latex apssamp.tex
%
\documentclass[%
%reprint,
superscriptaddress,
%groupedaddress,
%unsortedaddress,
%runinaddress,
%frontmatterverbose, 
%preprint,
%preprintnumbers,
%nofootinbib,
%nobibnotes,
%bibnotes,
 amsmath,amssymb,
%aps,
pra,
showkeys,
%prb,
%rmp,
%prstab,
%prstper,
%floatfix,
]{revtex4-2}


\usepackage{graphicx}% Include figure files
\usepackage{dcolumn}% Align table columns on decimal point
\usepackage{bm}% bold math
\usepackage{siunitx} % add this in your preamble
\usepackage[ruled,vlined]{algorithm2e}

\usepackage[colorlinks=true,linkcolor=blue,citecolor=blue,urlcolor=blue]{hyperref}
\usepackage[capitalize,noabbrev]{cleveref}
\usepackage{stmaryrd}

\sisetup{table-number-alignment = center, table-format = 2.2}
%%for legend%%%%%%%%%%%%%%%%%%%%%%%%%%%%%%%%%%%%%%%%%%%%%%%%%%%%%%%%%
\usepackage[table]{xcolor} %
\definecolor{plotblue}{HTML}{0000FF}
\definecolor{plotcyan}{HTML}{00BFBF}
\definecolor{plotred}{HTML}{FF0000}
\definecolor{plotgreen}{HTML}{008000}
\definecolor{plotgray}{HTML}{404040}



\newcommand{\colorline}[2][0.8em]{%
  \textcolor{#2}{\rule[0.5ex]{#1}{0.5pt}}%
}


\newcommand{\solidline}[2][0.9em]{%
  \textcolor{#2}{\rule[0.5ex]{#1}{0.6pt}}%
}

% Dotted line (M2)
\newcommand{\dottedline}[2][0.9em]{%
  \textcolor{#2}{\rule[0.5ex]{0.1em}{0.6pt}\hspace{0.25em}\rule[0.5ex]{0.1em}{0.6pt}\hspace{0.25em}\rule[0.5ex]{0.1em}{0.6pt}}%
}

% Dash-dotted line (M3)
\newcommand{\dashdotline}[2][0.9em]{%
  \textcolor{#2}{\rule[0.5ex]{0.4em}{0.6pt}\hspace{0.25em}\rule[0.5ex]{0.1em}{0.6pt}\hspace{0.25em}\rule[0.5ex]{0.4em}{0.6pt}}%
}

% Dashed line (M4)
\newcommand{\dashedline}[2][0.9em]{%
  \textcolor{#2}{\rule[0.5ex]{0.4em}{0.6pt}\hspace{0.25em}\rule[0.5ex]{0.4em}{0.6pt}}%
}

% Thick gray line (M5)
\newcommand{\thickgray}[2][0.9em]{%
  \textcolor{#2}{\rule[0.5ex]{#1}{1.2pt}}%
}
%%%%%%%%%%%%%%%%%%%%%%%%%%%%%%%%%%%%%%%%%%%%%%%%%%%%%%%%%%%%%%%%%%%%%


%\usepackage{hyperref}% add hypertext capabilities
%\usepackage[mathlines]{lineno}% Enable numbering of text and display math
%\linenumbers\relax % Commence numbering lines

%\usepackage[showframe,%Uncomment any one of the following lines to test 
%%scale=0.7, marginratio={1:1, 2:3}, ignoreall,% default settings
%%text={7in,10in},centering,
%%margin=1.5in,
%%total={6.5in,8.75in}, top=1.2in, left=0.9in, includefoot,
%%height=10in,a5paper,hmargin={3cm,0.8in},
%]{geometry}

\begin{document}

\preprint{APS/123-QED}

\title{A hybrid wall-modeled immersed boundary method with mixed subgrid-scale modeling for heterogeneous turbulent flows over complex terrain}% Force line breaks with \\
%\thanks{A footnote to the article title}%

\author{Arjun Ajay}
\email{arjunajay100@gmail.com}
\affiliation{School of Engineering, The University of British Columbia--Okanagan, 3333 University Way, Kelowna, British Columbia, Canada}

\author{Jagdeep Singh}
\thanks{Corresponding Author}
\email{jagdeep.singh@ubc.ca}

\affiliation{School of Engineering, The University of British Columbia--Okanagan, 3333 University Way, Kelowna, British Columbia, Canada}

\author{Sebastiano Stipa}
\affiliation{Environmental and Applied Fluid Dynamics Dept., von Karman Institute for Fluid Dynamics, Sint-Genesius-Rode, Belgium}

\author{Joshua Brinkerhoff}
\email{joshua.brinkerhoff@ubc.ca}
\affiliation{School of Engineering, The University of British Columbia--Okanagan, 3333 University Way, Kelowna, British Columbia, Canada}


\date{\today}% It is always \today, today,
             %  but any date may be explicitly specified

\begin{abstract}
	Large eddy simulation (LES) of atmospheric boundary layer (ABL) flow over complex terrain remains challenging due to the combined effects of steep topography, high Reynolds numbers, stratification, and unresolved turbulence. Accurate representation of surface geometry, near-wall processes, and subgrid-scale (SGS) dynamics is essential for predictive simulations, particularly for wind energy applications. In this study, we develop and assess a hybrid wall-modeled immersed boundary method (IBM) coupled with a mixed SGS turbulence closure for LES of ABL flow over complex terrain. The wall model is designed to operate on relatively coarse grids while still accurately representing near-wall stresses in separated and non-equilibrium flows characteristic of complex terrain. The SGS formulation combines a scale-similarity model with an anisotropic minimum-dissipation model to provide a balanced representation of energy transfer and to preserve the directional features of the unresolved scales. The proposed framework is used to simulate three LES studies: a flat terrain, an idealized three-dimensional cosine-squared hill, and a realistic natural-terrain case. First- and second-order flow statistics, as well as spectral characteristics, are evaluated against several widely used advanced SGS models, including the dynamic Smagorinsky model and the Lagrangian-averaged scale-invariant and scale-dependent dynamic formulations. The results demonstrate that the proposed hybrid wall-modeled IBM and mixed SGS model combination yield improved representation of mean flow and turbulence intensity over complex terrains compared to conventional approaches. These findings highlight the importance of jointly addressing terrain representation, wall modeling, and SGS closure to improve the reliability of LES for ABL flows over heterogeneous and topographically complex surfaces.
\end{abstract}

%\keywords{atmospheric boundary layer, large eddy simulation, complex terrain, immersed boundary method, wall modeling, subgrid-scale models, intermittency}

\maketitle

%\tableofcontents
\section{\label{sec:intro} Introduction}

Understanding the physical processes governing transport and exchange between the Earth’s surface and the atmosphere within the atmospheric boundary layer (ABL) is important for applications such as weather and climate prediction, air-quality assessment, and wind energy development \cite{chow2019crossing, lehner2018current, clifton2022research}. In regions of complex terrain, the evolution of the ABL state in space and time is strongly influenced by topographical features and land-surface heterogeneity, giving rise to flow separation, enhanced shear, and secondary circulations that complicate both physical interpretation and numerical modeling \cite{chow2019crossing, giani2024sensitivity}. These processes are particularly relevant for wind energy applications, where terrain-induced flow heterogeneity and wake/terrain interactions directly affect turbine performance \cite{lehner2018current}.

Numerical weather prediction (NWP) models such as WRF \cite{michalakes2001, michalakes2005weather} are widely used to simulate atmospheric flow over complex terrain. These models typically resolve mesoscale dynamics at horizontal grid spacings of order $\mathcal{O}(1)$ - $\mathcal{O}(10)\,\mathrm{km}$ \citep{giani2024sensitivity,powers2017weather}. However, capturing terrain-induced flow features and land/atmosphere interactions requires much finer spatial resolution, often on the order of $\mathcal{O}(10)\,\mathrm{m}$. Notably, such resolution is also required for simulation where wind turbines are represented using actuator models \cite{wu2013simulation, maas2022wake, cheung2023investigations, lanzilao2024parametric, stipaafm2024}. At these grid spacings, the use of terrain-following coordinates can lead to grid skewness and spurious numerical errors when applied over steep or irregular topography, which degrade solution accuracy and may compromise numerical stability \cite{lundquist2010immersed, lundquist2012immersed, daniels2016new, arthur2019using}. To address these limitations, multiscale nested WRF-LES modeling frameworks are increasingly adopted, in which mesoscale simulations provide dynamically consistent boundary conditions to embedded LES domains at much finer resolution \cite{mirocha2010implementation,talbot2012nested,munoz2014bridging,munoz2015stochastic,powers2017weather,liu2020simulation}. However, turbulence modeling is particularly challenging in nested WRF-LES frameworks, as the spatial resolution must jump over the \textit{gray-zone}, where the grid size is too coarse for an LES and too fine for planetary boundary layer parameterizations \cite{wyngaard2004toward, chow2019crossing}.

\citet{chow2019crossing} identified three major challenges associated with numerical modeling of atmospheric flow in mountainous regions. The first concerns terrain representation, as steep slopes often challenge the ability of numerical grids to capture surface features accurately. Another difficulty arises from turbulence modeling, since complex landscapes strongly affect mixing processes and the dissipation of kinetic energy, requiring careful treatment of turbulence parameterizations. A third challenge involves the effects of atmospheric stratification, because mountain-driven circulations, surface cooling and heating can significantly alter local temperature and moisture distributions, complicating the representation of convective processes, particularly in coarse-resolution models.

With advances in computational power and the development of meso-to-microscale coupling strategies~\citep{allaerts2020development,Ajay2025ProfileAssimilation}, large eddy simulation (LES) has become increasingly feasible for representing realistic atmospheric flow, capturing land/atmosphere exchanges and transport processes, incorporating spatial heterogeneity, and assessing wind farm performance and layout optimization. However, the challenges identified by \citet{chow2019crossing} for NWP remain relevant in LES, albeit in different forms. Achieving a horizontal resolution of $\mathcal{O}(10\,\mathrm{m})$ is still computationally demanding when simulating atmospheric boundary layer flow over wind farms sited on complex terrain, where flow separation, wake development, and land/atmospheric exchange processes must be resolved and unresolved small-scale flow physics must be accurately modeled via subgrid scale (SGS) models. Similarly, representing steep topography continues to require robust methods to avoid numerical artifacts associated with terrain-following grids \cite{lundquist2012immersed,bhuiyan2022scale,chow2019crossing}.  
When the grid resolution only partially resolves the turbulent motions, the representation of steep and irregular topography becomes critical and often requires advanced techniques such as the immersed boundary method (IBM) to accurately capture surface-flow interactions \cite{lundquist2012immersed,bao2018large,giani2024sensitivity}. Under these conditions, the influence of terrain and surface heterogeneity places additional demands on SGS modeling, which must recover the unresolved turbulence kinetic energy (TKE) generated by terrain-induced motions that cannot be explicitly resolved at the grid scale \cite{bhuiyan2022scale,singh2024impact}. The interaction between stratification and turbulence over complex terrain adds further complexity, as stability strongly influences flow separation, mixing, and energy transport. Addressing these interconnected issues is essential for improving the predictive capability of LES and advancing its application to wind-farm flows and land-atmosphere exchange over complex surfaces. 
%Furthermore, the interaction of stratification with turbulence in such environments places stringent demands on SGS modeling.  
%The goal of this paper is to address these fundamental challenges in the context of LES of ABL over complex terrain.

Accurate geometrical representation of complex terrain has long been a limiting factor in numerical modeling of ABL flow \cite{bao2018large,giani2024sensitivity}. Numerical models often struggle to capture ABL flow over complex terrain due to inherent nonlinearity and non-equilibrium effects \cite{uchida2003large,wan2007evaluation,cao2012numerical,yang2021study}.
Traditional approaches based on terrain-following coordinates (TFC) \cite{clark1977small,schar2002new} are susceptible to significant numerical errors when applied to steep slopes, leading to grid distortion, spurious pressure gradients, and a reduction in overall solution accuracy \cite{giani2024sensitivity}. To overcome these limitations, IBM has been developed to represent complex geometries on Cartesian grids without the need for body-fitted coordinates. Pioneered by \citet{peskin1972flow} for biological flow problems, IBM has since become one of the most widely used approaches for representing complex geometries across a broad range of applications, including urban environments and mountainous terrain \cite{tseng2006modeling,diebold2013flow,bou2009effects,lundquist2010immersed,lundquist2012immersed,bao2018large,alam2018large}. A key advantage of IBM lies in its ability to retain a regular Cartesian computational mesh while accommodating geometrically complex boundaries. Several formulations of IBM exist, differing primarily in how the influence of the boundary is incorporated into the governing equations; comprehensive reviews are provided by \citet{iaccarino2003immersed}, and \citet{mittal2005immersed}. Early formulations relied on distributed body-force representations, in which the effect of the boundary was spread over multiple grid cells \cite{Peskin2002}, whereas more recent sharp-interface methods enforce boundary conditions locally at the fluid–solid interface \cite{mohd1997combined,fadlun2000combined,gilmanov2005hybrid}. Within the sharp-interface class, approaches include cut-cell methods \cite{Meyer2010,Kirkpatrick2003}, which modify control volumes to conform to the embedded geometry, and immersed-interface or reconstruction-based methods, which impose boundary conditions through local interpolation or reconstruction procedures without altering the mesh topology. A widely used subset of the latter is the direct-forcing immersed boundary approach \cite{Deleon2018,gilmanov2005hybrid}, in which flow variables in boundary-adjacent cells are reconstructed directly from prescribed boundary conditions and nearby fluid values, eliminating the need for explicit evaluation of body-force terms. This approach, adopted in the present study, significantly simplifies grid generation and avoids the mesh distortions commonly associated with TFC systems, particularly over steep or irregular surfaces.

In LES of ABL flows, resolving the near-wall region of an immersed-body through grid refinement is generally infeasible. Although wall-resolved simulations may employ linear velocity reconstruction at the immersed surface when the viscous sub-layer is resolved \cite{gilmanov2005hybrid,ge2007numerical,borazjani2008curvilinear}, such resolutions are prohibitively expensive for the high Reynolds number regimes typical of atmospheric boundary layers. As a result, a wall-model must be applied at the immersed boundary interface, independent of the underlying terrain characteristics. Wall modeling within the IBM framework is notably more challenging than in terrain-following coordinate (TFC) approaches because the immersed surface seldom coincides with grid points. In practice, this misalignment is commonly handled using smearing-based treatments \cite{diebold2013flow,fang2016intercomparison} or linear interpolation schemes to impose boundary conditions exactly at the immersed surface. These approaches can introduce non-negligible errors, particularly in the prediction of near-wall stresses, motivating the development of improved wall-model formulations. Existing wall-modeled immersed boundary approaches are broadly classified into velocity-reconstruction and shear-stress-reconstruction methods \cite{chester2007modeling,liu2020wall}. However, many current methods struggle to accurately capture mean velocity and velocity variance profiles in highly heterogeneous turbulent flows. Such methods often require excessive grid refinement to represent non-equilibrium effects, even at moderate Reynolds numbers \cite{diebold2013flow,liu2020wall}. %The present work addresses these limitations through the development of a hybrid wall-modeled immersed boundary method specifically tailored for high Reynolds-number ABL flows and suitable for use on relatively coarser grids.

In addition to wall-modeled immersed boundary methods, the treatment of unresolved turbulence remains a central challenge in LES, especially over complex surfaces where topography significantly alters turbulence generation, anisotropy, and dissipation. Classical eddy-viscosity models, such as the Smagorinsky model \citep{smagorinsky1963general}, have been widely used but tend to be overly-dissipative. Dynamic models \citep{germano1991dynamic} improve this limitation by allowing the model coefficient to adapt to local flow conditions, while scale-dependent variants \citep{porte2000scale,bou2005scale} further enhance accuracy in stratified or shear-dominated flows. Mixed models, which combine structural and dissipative components, offer a more balanced representation of energy transfer between resolved and subgrid scales and are able to preserve the directional features of SGS stresses \cite{bardina1983improved,vreman1994formulation,horiuti1997new}. Although these models often yield improved accuracy, they typically require multiple filtering operations, thereby increasing computational cost.

When grid resolution only partially resolves turbulent motions, the treatment of unresolved turbulence becomes equally critical to the representation of complex terrain. Most LES formulations rely on eddy-viscosity-based SGS models, which assume that unresolved motions are locally isotropic and that subgrid-scale dissipation can be represented through a mean, equilibrium process \cite{smagorinsky1963general,sagaut2006large}. These assumptions imply alignment of the eigenframe of the SGS stress tensor with that of the resolved-scale strain-rate tensor and form the basis of many commonly used closures. 

Atmospheric flow over a three-dimensional hill provides a natural test case for examining these issues. Terrain-induced pressure gradients, curvature, and separation introduce strong spatial variability, while interactions between SGS closures and wall-modeled immersed boundary treatments directly influence near-surface turbulence. Differences in how SGS models represent or suppress intermittent fluctuations can therefore be used to identify limitations in their underlying assumptions. Among the various SGS models used for LES of ABL flow over complex terrain, multiple studies have reported poorer performance of dynamic SGS models, including the dynamic Smagorinsky model \cite{ghosal1995dynamic} and the dynamic Lagrangian averaged scale-invariant model (LASI) \cite{meneveau2000scale}, compared to the standard Smagorinsky model \cite{smagorinsky1963general}. This behavior has been attributed to underestimation of the model coefficient near the surface \cite{wan2007evaluation,iizuka2004performance,yang2020study}. To address this limitation, Wan et al.~\cite{wan2007evaluation} advocated the scale-dependent dynamic Lagrangian model (LASD) \cite{bou2005scale}, which relaxes the assumption of scale invariance for the model coefficient used in the Germano dynamic SGS framework.

%When grid resolution is insufficient to resolve energetic eddies, particularly near steep and heterogeneous terrain, the assumption of small-scale isotropy becomes increasingly questionable and places a greater burden on SGS models. Under such strongly heterogeneous conditions, SGS closures must recover unresolved turbulence kinetic energy (TKE) and represent partial topographic effects, tasks that remain challenging and require continued refinement. Among the various SGS models used for LES of ABL flow over complex terrain, multiple studies indicate poorer performance of dynamic subgrid-scale models, including the dynamic Smagorinsky model \cite{ghosal1995dynamic} and the dynamic Lagrangian averaged scale-invariant model (LASI) \cite{meneveau2000scale}, compared to the standard Smagorinsky model \cite{smagorinsky1963}. This behavior is attributed to underestimation of the model coefficient near the surface \cite{wan2007evaluation,iizuka2004performance,yang2020study}. Wan et al. \cite{wan2007evaluation} advocate the use of the scale-dependent dynamic Lagrangian model (LASD) \cite{bou2005scale}, which removes the assumption of scale invariance for the model coefficient used in the Germano dynamic subgrid-scale model (DSM) \cite{ghosal1995dynamic}.

In the present study, we develop a hybrid wall-modeled immersed boundary method that combines shear-stress and velocity reconstruction methods for application to atmospheric boundary layer (ABL) flows. This IBM formulation is coupled with a mixed subgrid-scale (SGS) model that employs a scale-similarity model \cite{bardina1983improved} as the structural component and an anisotropic minimum-dissipation model as the functional component \cite{rozema2015minimum,verstappen2018much}. This mixed SGS model was evaluated in our recent work, where it shown improved accuracy as well as higher computational efficiency than other mixed formulations \cite{ajay2025mixed}. In the present study, the performance of this approach and that of several other advanced eddy-viscosity models --- including LASI, LASD and localized dynamic Smagorinsky models --- is examined in combination with the proposed hybrid wall-modeled IBM formulation for idealized complex terrain using first- and second-order turbulence statistics and energy-spectral analyses. Finally, the best configuration is applied to a realistic terrain case based on the Bolund Hill site. While stratification effects are an important aspect of atmospheric flows over complex terrain, as described above, the present study focuses on neutral conditions and does not explicitly consider thermally stratified regimes.

The remainder of the paper is organized as follows. Section~\ref{sec:method} describes the methodological framework employed in this study. Section~\ref{sec:methodles} introduces the large eddy simulation formulation, including the governing equations and numerical setup. Section~\ref{sec:methodturbulencemodel} presents the turbulence closures considered, with emphasis on the mixed subgrid-scale model and reference eddy-viscosity models. Section~\ref{sec:methodibm} details the hybrid wall-modeled immersed boundary method used to represent complex terrain at high Reynolds numbers. Section~\ref{sec:results} presents the results and discusses the performance of the wall model and subgrid-scale closures over complex terrain. Finally, Section~\ref{sec:conclusion} summarizes the main findings, limitations and provides concluding remarks.

\section{\label{sec:method} Methods}
\subsection{\label{sec:methodles} Large eddy simulation}
The present study employs the Toolbox fOr Stratified Convective Atmospheres (TOSCA), an open-source large eddy simulation (LES) framework based on a finite-volume formulation in generalized curvilinear coordinates (GCC). TOSCA has been extensively validated across a wide range of atmospheric boundary layer and wind energy applications, including wind-farm wake interactions, gravity-wave interactions with wind farms and heterogeneous environment, and meso-to-microscale (MMC) coupling \cite{stipa2024tosca,stipaafm2024,ajay2025mixed,Ajay2025ProfileAssimilation}. Governing equations are obtained by applying a spatial low-pass filter to the incompressible Navier-Stokes equations, thereby separating the resolved large-scale motions from the unresolved subgrid-scale contributions \cite{sagaut2006large}. The characteristic filter width is defined as $\Delta_{\hbox{\tiny LES}}=\sqrt[3]{\Delta_x\Delta_y\Delta_z}$, based on the local grid spacing in the streamwise, spanwise, and vertical directions. The filtered continuity and momentum equations are given by
\begin{equation}
	\frac{\partial \overline{u}_{i}}{\partial x_i} = 0,
	\label{eq:nseLES1}
\end{equation}
%
\begin{equation}
	\frac{\partial \overline{u}_i}{\partial t}
	+ \frac{\partial \left(\overline{u}_{i}\overline{u}_{j}\right)}{\partial x_j}
	= -\frac{\partial \overline{p}}{\partial x_i}
	+ \frac{\partial}{\partial x_j}
	\left[\nu\left(\frac{\partial \overline{u}_i}{\partial x_j}
	+ \frac{\partial \overline{u}_j}{\partial x_i}\right)\right]
	- \frac{\partial \tau_{ij}}{\partial x_j} + f^{IB}_i.
	\label{eq:nseLES2}
\end{equation}
%
Here, the overbar denotes spatially filtered quantities, $\overline{u}_i$ represents the resolved velocity component, $\overline{p}$ is the pressure normalized by the constant fluid density, and $\nu$ is the kinematic viscosity. The last term $f^{IB}_i$ in Eq.~\eqref{eq:nseLES2} represents the body force at the interface of the immersed-body. Within the class of the direct forcing IBM, this term imposes a boundary constraint on the velocity field. The stress term $\tau_{ij}=\overline{u_i u_j}-\overline{u}_i\overline{u}_j$ corresponds to the subgrid-scale (SGS) stress tensor and represents the influence of unresolved motions on the resolved flow field \cite{ferziger1977large,lesieur1996new,mason1994large,rogallo1984numerical}. Since $\overline{u_i u_j}\neq\overline{u}_i\overline{u}_j$, a closure problem arises. To close the system of Eqs.~\ref{eq:nseLES1} and~\ref{eq:nseLES2}, a turbulence closure model is therefore introduced to approximate the SGS stress term. The resulting equations, together with appropriate boundary conditions, are discretized and solved to obtain an approximation of the resolved velocity field $\overline{\bm{u}}$. A wide range of SGS models exists and they are generally classified into three main categories: (i) functional, (ii) structural, and (iii) mixed models. In the present study, we employ the mixed Bardina-anisotropic minimum dissipation (BAMD) model, which combines a structural scale-similarity formulation \cite{bardina1983improved} with an anisotropic minimum-dissipation eddy-viscosity closure \cite{rozema2015minimum,verstappen2018much,vreugdenhil2018large}. The BAMD model has been recently validated for heterogeneous and highly inhomogeneous turbulent flow conditions, including complex terrain and atmospheric boundary layer applications \cite{ajay2025mixed}. By blending functional and structural SGS representations, the BAMD model provides both numerical robustness and improved representation of scale interactions. Details of the BAMD formulation and its implementation are described in the following section.

\subsection{\label{sec:methodturbulencemodel}Turbulence closures}
Accurate LES requires both faithful structural representation of the unresolved motions and sufficient dissipation of SGS energy. These objectives are commonly achieved through mixed models, in which a functional eddy-viscosity closure is combined with a structural component. The earliest mixed formulation was proposed by \citet{bardina1983improved}, who coupled the Smagorinsky model~\cite{smagorinsky1963} with a scale-similarity closure. In generalized curvilinear coordinates, the SGS stress contribution to the momentum equations is written as
\begin{equation}
	T_i^{r} = {T_i^{r}}^{\alpha} + {T_i^{r}}^{\beta},
	\label{eq:contributionLES2}
\end{equation}
where
\begin{equation}
	T_i^{r}
	= \frac{1}{J}\frac{\partial l^{r}}{\partial x_j}\tau_{ij}
	= \frac{1}{J}\left( \overline{U^{r}u_i} - \overline{U}^{r}\,\overline{u}_i \right),
	\label{eq:recastCurv}
\end{equation}
with $J$ denoting the Jacobian of the coordinate transformation. The functional contribution is further decomposed into deviatoric and isotropic parts,
\begin{equation}
	{T_i^{r}}^{\alpha}
	= {T_i^{r}}^{\alpha,\hbox{\tiny dev}}
	+ {T_i^{r}}^{\alpha,\hbox{\tiny iso}},
	\label{eq:contributionLESalpha}
\end{equation}
where the deviatoric term is modeled using an eddy-viscosity assumption,
\begin{equation}
	{T_i^{r}}^{\alpha,\hbox{\tiny dev}}
	= -2\nu_{\tau}\,\overline{S}_i^{r}.
	\label{eq:isoparteddy}
\end{equation}
The structural contribution is represented using the scale-similarity model,
\begin{equation}
	{T_i^{r}}^{\beta}
	= C_b \left(
	\widetilde{\overline{U^{r}}\overline{u}_i}
	- \widetilde{\overline{U^{r}}}\,\widetilde{\overline{u}}_i
	\right),
	\label{eq:bardina}
\end{equation}
where $\overline{U}^{r}$ denotes the filtered contravariant velocity components. The contravariant and Cartesian velocities are related through
\begin{equation}
	\overline{U}^{r}
	= \frac{\partial l^{r}}{\partial x_i}\,\overline{u}_i,
	\qquad
	\overline{u}_i
	= \frac{\partial x_i}{\partial l^{r}}\,\overline{U}^{r}.
\end{equation}
The Bardina model constant is set to $C_b=1.0$, which ensures Galilean invariance as demonstrated by \citet{speziale1985galilean}. The primary filter $(\overline{\cdot})$ corresponds to the grid filter of width $\Delta_{\hbox{\tiny LES}}$, while the test filter $(\widetilde{\cdot})$ has a width $2\Delta_{\hbox{\tiny LES}}$.

For the functional component, the anisotropic minimum dissipation (AMD) model is adopted due to its low-dissipation characteristics and its suitability for anisotropic grids. Atmospheric boundary layer simulations over complex terrain typically employ grids with strong vertical refinement, whereas many conventional eddy-viscosity models assume isotropic resolution. The AMD model explicitly accounts for grid anisotropy and has been shown to improve performance in such configurations \cite{rozema2015minimum,moser2021statistical}.

To ensure frame invariance on anisotropic grids, the AMD formulation employs scaled velocity and spatial coordinates \cite{verstappen2018much}. The eddy viscosity is given by
\begin{equation}
	\nu_t
	= -\left(C\Delta^2\right)
	\frac{
		\left(\hat{\partial}_k \hat{u}_i\right)
		\left(\hat{\partial}_k \hat{u}_j\right)
		\hat{S}_{ij}
	}{
		\left(\hat{\partial}_l \hat{u}_m\right)
		\left(\hat{\partial}_l \hat{u}_m\right)
	},
	\label{eq:subgridviscosity}
\end{equation}
where $C$ is the modified Poincar\'e constant. The effective filter width $\Delta$ is computed as the square root of the harmonic mean of the directional filter widths,
\begin{equation}
	\frac{1}{\Delta^2}
	= \frac{1}{3}
	\left(
	\frac{1}{\Delta_x^2}
	+ \frac{1}{\Delta_y^2}
	+ \frac{1}{\Delta_z^2}
	\right),
\end{equation}
where $\Delta_x$, $\Delta_y$, and $\Delta_z$ denote the local grid spacings in the Cartesian directions. The modified Poincar\'e constant is set to $C=0.12$ for the mixed modeling approach. This value is chosen based on sensitivity analyses and the numerical discretization employed, which uses a fourth--order central scheme with hyperviscous stabilization \cite{ajay2025mixed}. The scaled variables are defined as
\[
\hat{x}_i = \frac{x_i}{\Delta_i},
\qquad
\hat{u}_i = \frac{\overline{u}_i}{\Delta_i},
\qquad
\hat{\partial}_i \hat{u}_j
= \frac{\Delta_i}{\Delta_j}\,\partial_i \overline{u}_j,
\]
and the scaled rate-of-strain tensor is
\[
\hat{S}_{ij}
= \frac{1}{2}\left(
\hat{\partial}_i \hat{u}_j
+ \hat{\partial}_j \hat{u}_i
\right).
\]
Repeated indices in Eq.~\eqref{eq:subgridviscosity} imply summation over spatial dimensions, ensuring that $\nu_t$ is a scalar quantity and application of Eq. \eqref{eq:subgridviscosity} in curvilinear coordinates involves computation of the velocity gradients in curvilinear coordinates followed by transformation to the Cartesian gradients using the metric terms as $\frac{\partial \overline{u}_i}{\partial x_j} = \frac{\partial \overline{u}_i}{\partial l^r} \frac{\partial l^r}{\partial x_j}$.

Finally, the deviatoric component of the functional closure, given by Eq.~\eqref{eq:isoparteddy}-\eqref{eq:subgridviscosity}, together with the structural contribution defined by Eq.~\eqref{eq:bardina}, are substituted into Eq.~\eqref{eq:contributionLES2} to parameterize $T_i^{r}$. In the present study, we also consider other advanced eddy-viscosity models, including the Lagrangian averaged scale-dependent (LASD) model \cite{bou2005scale}, the Lagrangian averaged scale-invariant (LASI) model \cite{meneveau2000scale}, and the localized dynamic Smagorinsky (DSM) model \cite{ghosal1995dynamic}. For brevity, the formulation of these models in generalized curvilinear coordinates is not presented here; interested readers are referred to \citet{ajay2025mixed} for detailed descriptions and implementations. %In the following section, we describe the newly developed hybrid wall-modeled immersed boundary method.

\subsection{\label{sec:methodibm} Hybrid wall-modeled immersed-boundary method} 

This study employs a direct-forcing immersed boundary method, in which the body-force term $f^{IB}$ appearing in Eq.~\ref{eq:nseLES2} is not evaluated explicitly. Instead, the influence of the immersed boundary is imposed implicitly through reconstruction of the velocity field at the immersed-body interface. The immersed-body interface is identified using a ray-tracing algorithm \cite{borazjani2008curvilinear}, with the immersed-body surface discretized using an unstructured triangular mesh. 

For each computational grid cell, a ray is cast in an arbitrary direction from the cell center, and the number of intersections with the triangular elements of the immersed-body surface determines whether the cell center lies inside or outside the body. An odd number of intersections indicates that the cell center lies inside the solid, whereas an even number indicates that it lies outside. Based on this classification, the computational grid is partitioned into three categories: fluid cells located outside the immersed-body, solid cells located within the body, and immersed-body interface cells, defined as the first layer of solid cells whose centers lie inside the immersed-body. Additionally, for each immersed-body interface cell, the closest triangular element of the immersed-body surface is identified, which is used to compute the unit normal direction to the IBM surface for the interface cell.

In wall-resolved simulations, where the grid spacing is sufficiently fine to resolve the viscous sublayer, a linear reconstruction of the velocity at the immersed-body interface cells is generally adequate \cite{gilmanov2005hybrid}. For the high-Reynolds-number flows typical of atmospheric boundary layer LES, however, such near-wall resolutions are computationally prohibitive. As a result, a wall model must be applied at the immersed-body interface to represent the unresolved near-wall dynamics. The existing literature on wall-modeled IBM can be broadly classified into two approaches: velocity reconstruction methods \cite{Choi2007,Deleon2018} and shear-stress reconstruction methods \cite{chester2007modeling,liu2020wall}. Velocity reconstruction methods employ a wall-model formulation—commonly based on the logarithmic law of the wall—to estimate the velocity at the first off-wall point and then reconstruct values at immersed-body interface cells through interpolation or extrapolation. Shear-stress reconstruction methods, in contrast, compute the wall shear stress directly from a wall model and impose it as a boundary condition at the immersed-body interface. The modeled stress is applied through forcing terms in the momentum equations.

The present work adopts a hybrid strategy that incorporates elements of both approaches. The wall shear stress at the immersed-body interface is obtained from the tangential fluid velocity sampled from fluid cells near the interface, as described below. In parallel, the ghost-node velocity boundary condition at the immersed-body interface cells is obtained by extrapolating the velocity field from sampling points within the fluid region. This combined treatment improves the accuracy of the IBM wall model by using shear-stress reconstruction to correctly impose surface drag at the immersed-body interface, while simultaneously providing consistent velocity values for the gradient stencils associated with the hybrid staggered/non-staggered discretization employed in TOSCA \cite{ajay2025mixed}. A schematic illustration of the method is provided in Fig.~\ref{fig:IBMwallmodel}.


\begin{figure}[htb!]
    \centering
    \includegraphics[width=0.5\linewidth]{plots/ibmfig.pdf}
\caption{Schematic illustration of the hybrid wall-modeled immersed boundary method. Blue solid circles denote solid cells, green circles indicate immersed-body interface cells (IBM interface cells), and red solid circles represent fluid cells. The purple curve represents the immersed-body surface, while the orange dashed lines mark the faces of the immersed-body interface cells where the wall shear stress is applied. For each interface face, the wall-normal direction $\hat{\boldsymbol{n}}$ is defined as the unit normal to the closest immersed-body surface triangular element associated with the IBM interface cell, and distances $\delta$ are measured along this direction. Additionally, the ghost cell (immersed-body interface cell) $G$, located inside the solid region, is used to impose boundary conditions through velocity reconstruction.}
    \label{fig:IBMwallmodel}
\end{figure}

\subsubsection{Shear-stress reconstruction}

Shear stress is reconstructed at faces shared by an immersed-body interface cell and a neighboring fluid cell. Let $\boldsymbol{\mathsf{n}}$ denote the unit normal to the closest immersed-body surface triangular element associated with the IBM interface cell, and let $\boldsymbol{\mathsf{u}}$ denote the fluid velocity sampled at a distance $\delta$ from the surface along $\boldsymbol{\mathsf{n}}$. For isotropic grids, $\delta=\Delta$, where $\Delta$ is the grid spacing, whereas for anisotropic grids, $\delta$ represents the approximate local grid spacing in the wall-normal direction.

The tangential component of the sampled velocity is obtained using the projection operator,
\[
\mathbb{P}_t = \mathbb{I} - \boldsymbol{\mathsf{n}} \otimes \boldsymbol{\mathsf{n}},
\qquad
\boldsymbol{\mathsf{u}}_t = \mathbb{P}_t\,\boldsymbol{\mathsf{u}},
\qquad
u_t = \lVert \boldsymbol{\mathsf{u}}_t \rVert .
\]
Following \citet{chester2007modeling}, a local orthonormal coordinate frame aligned with the immersed-body surface is defined as
\[
\boldsymbol{\varepsilon}_1 = \frac{\boldsymbol{\mathsf{u}}_t}{u_t},
\qquad
\boldsymbol{\varepsilon}_2 = \boldsymbol{\mathsf{n}} \times \boldsymbol{\varepsilon}_1,
\qquad
\boldsymbol{\varepsilon}_3 = \boldsymbol{\mathsf{n}},
\]
with the corresponding rotation matrix
\[
\mathbf{Q}
=
\big[\,\boldsymbol{\varepsilon}_1\;\boldsymbol{\varepsilon}_2\;\boldsymbol{\varepsilon}_3\,\big],
\]
which maps quantities between the global and local wall-parallel coordinate systems.

Assuming that the tangential velocity follows a rough-wall logarithmic profile, the magnitude of the wall shear stress is given by
\begin{equation}
\tau_w
=
\left[
\frac{\kappa\,u_t}{\ln\!\left(\delta/z_0\right)}
\right]^2 ,
\label{eq:IBM_tau_final}
\end{equation}
where $\kappa = 0.4$ is the von K\'arm\'an constant and $z_0$ denotes the aerodynamic roughness length.
Although the face center $x_c$ does not, in general, lie on the immersed-body surface, the quantity $\tau_w$ is taken to represent the $\mathcal{T}_{13} = \mathcal{T}_{31}$ component of the shear stress tensor in the local coordinate system, with all other components set to zero. Accordingly, the local shear stress tensor $\mathcal{T}^{(\mathrm{loc})}$ is expressed as
\begin{equation}
\mathcal{T}^{(\mathrm{loc})}
=
\tau_w
\big(
\boldsymbol{\varepsilon}_1 \otimes \boldsymbol{\varepsilon}_3
+
\boldsymbol{\varepsilon}_3 \otimes \boldsymbol{\varepsilon}_1
\big).
\end{equation}
The corresponding stress tensor imposed at the immersed-body interface face in the global coordinate system
is obtained through the transformation
\begin{equation}
\mathcal{T}
=
\mathbf{Q}\,
\mathcal{T}^{(\mathrm{loc})}\,
\mathbf{Q}^{\!\top}.
\end{equation}
Consistent with \citet{liu2020wall}, this stress is applied only at immersed-boundary interface faces and is set to zero inside the immersed-body.

\subsubsection{Velocity reconstruction}
The second component of the hybrid wall model concerns the reconstruction of the velocity in the immersed-body interface cells which also act as ghost cells for the adjacent fluid cells. Two sampling points are defined along the wall-normal direction from the ghost-cell center at distances $0.5\delta$ and $1.5\delta$ from the immersed-body surface, denoted as points A and B in Fig.~\ref{fig:IBMwallmodel}. The velocity at point B, $\boldsymbol{\mathsf{u}}_B$, is obtained by trilinear interpolation from the background grid, and its tangential component is computed as
\begin{equation}
\boldsymbol{\mathsf{u}}_{B,t}
=
\mathbb{P}_t\,\boldsymbol{\mathsf{u}}_B .
\end{equation}

Assuming a logarithmic variation of the tangential velocity between the two sampling locations, the magnitude of the tangential velocity at point A is given by
\begin{equation}
u_{A,t}
=
u_{B,t}\,
\frac{\ln(0.5\delta/z_0)}{\ln(1.5\delta/z_0)},
\end{equation}
with its direction taken to be identical to that of $\boldsymbol{\mathsf{u}}_{B,t}$, where $z_0$ denotes the surface roughness length. Linear extrapolation of the tangential velocities from points A and B is then used to obtain the tangential component of the ghost-cell velocity, $\boldsymbol{\mathsf{u}}_{g,t}$. The wall-normal component of the ghost-cell velocity is determined by enforcing the no-penetration condition at the immersed-body surface, yielding the reconstructed ghost-cell velocity $\boldsymbol{\mathsf{u}}_g$. This reconstructed velocity is imposed as the boundary condition at the immersed-body interface and provides the required stencil values for evaluating velocity gradients in the vicinity of the interface.

\section{\label{sec:results}Results}
%\subsection{Representation of complex terrain}
In this section, we evaluate the performance of the hybrid wall modeled immersed boundary in capturing the aggregate impact of terrain on the flow for relatively coarse grids. For validation, we compare three representative cases: a flat terrain, a steep three-dimensional hill, and a realistic terrain featuring the steep Bolund hill. For clarity, a temporal average of quantity $Q$ is denoted as $\langle Q \rangle$, and a spatial and temporal average of quantity $Q$ is denoted as $\langle \!\langle Q \rangle \!\rangle $.

\subsection{Flat terrain}
The performance of the proposed hybrid wall model immersed boundary method is evaluated by simulating atmospheric boundary layer flow over flat terrain. Four cases are considered in which the lower wall is located at $z_w/dz = 1,\ 1.25,\ 1.5,\ 1.75$, respectively, where $dz$ is the grid size in the vertical direction, as illustrated in Fig.~\ref{fig:schematicFlat} for Cases 1 to 4. These simulations, performed using the hybrid wall model immersed boundary method, are compared against a baseline reference case that employs a conventional approach for modeling the solid surface (see, case-ref in Fig. \ref{fig:schematicFlat}a). Among the selected cases, Case 2 ($z_w/dz = 1.25$) and Case 4 ($z_w/dz = 1.75$) involve wall positions that are not coplanar with a face of the computational grid, allowing an assessment of the method’s accuracy and robustness in handling arbitrarily positioned solid boundaries.

A computational domain of size $L_x \times L_y \times L_z = 8 \text{ m} \times 4 \text{ m} \times 1 \text{ m}$ is considered, discretized into $N_x \times N_y \times N_z = 128 \times 64 \times 64$ finite volume cells. The domain height is set to $1.0~\mathrm{m}$, and the effective roughness length is specified as $z_0/L_z = 5.6 \times 10^{-5}$. A pressure controller \cite{stipa2024tosca} is used to drive the ABL flow over the flat terrain, with a velocity of $4~\mathrm{m/s}$ imposed at a height of $z = 0.1~\mathrm{m}$. The present computational configuration is closely aligned with the numerical setup of flat terrain reported by \cite{liu2020wall}.

\begin{figure}[htb!]
    \centering
    \includegraphics[width=0.65\linewidth]{plots/flat/flatterraincases.pdf}
    \caption{Position of the immersed flat surface relative to the grid for Cases~1--4, corresponding to
$z_w/dz = 1,\ 1.25,\ 1.5,\ \text{and}\ 1.75$, compared with the reference case without an
IBM surface (Case--Ref).
}
    \label{fig:schematicFlat}
\end{figure}

Fig.~\ref{fig:loglawFlat} presents a comparison between the theoretical logarithmic law and the simulated time- and planar-averaged streamwise velocity profiles. The wall-modeled IBM cases show excellent agreement with both the reference case and the theoretical profile. The only exception is Case~2, which slightly deviates from the expected trend. This discrepancy may arise from two primary factors. First, the wall location at $1.25\,dz $ is too close to the surface and lies outside the logarithmic region, which limits its consistency with the log-law behavior. Second, the wall in this case is not aligned with the computational grid. Consequently, the boundary conditions, which are defined at discrete grid points, must be imposed at intermediate locations within the fluid or the wall. This introduces approximations that may contribute to the deviation observed in Case~2.
\begin{figure}[htb!]
    \centering
    \begin{tabular}{c}
            \includegraphics[width=0.5\linewidth]{plots/flat/loglaw.pdf} \\
    \end{tabular}
\caption{Planar and time-averaged streamwise velocity profile for simulation over flat terrain, where terrain is represented using a hybrid wall-modeled approach.}
    \label{fig:loglawFlat}
\end{figure}

The variances of the streamwise, spanwise, and vertical velocity components are examined to evaluate the representation of turbulence. These variances are computed from the diagonal components of the resolved Reynolds stress tensor, $R_{ij} = \langle\!\langle {u_i' u_j'} \rangle\!\rangle
$, with the velocity fluctuations defined as $\sigma_{ii} = \sqrt{R_{ii}}$ and normalized by the friction velocity $u_*$. Figure~\ref{fig:sigmaUflat} shows the vertical profiles of the normalized streamwise ($\sigma_u$), spanwise ($\sigma_v$), and vertical ($\sigma_w$) velocity fluctuations. The wall-modeled immersed boundary method (IBM) cases closely follow the reference solution across all velocity components, indicating consistent reproduction of turbulence intensities.


\begin{figure}[htb!]
    \centering
    (a)
    \includegraphics[width=0.4\linewidth]{plots/flat/sigmaU.pdf} \\
    \begin{tabular}{cc}
            (b) & (c) \\
               \includegraphics[width=0.4\linewidth]{plots/flat/sigmaV.pdf} &
                  \includegraphics[width=0.4\linewidth]{plots/flat/sigmaW.pdf} \\
    \end{tabular}
\caption{Vertical profiles of the (a) streamwise, (b) spanwise, and (c) vertical components of the velocity variance computed from the resolved Reynolds stress tensor for a simulation over flat terrain.}
    \label{fig:sigmaUflat}
\end{figure}

\subsection{A steep three-dimensional hill}

To further evaluate the performance of the proposed framework in a more complex setting, the present LES results are compared with wind-tunnel measurements~\cite{cao2006experimental} and simulations conducted using a pseudo-spectral code~\cite{liu2020wall} for flow over a three-dimensional hill. The hill geometry follows a cosine-squared profile defined as
\begin{equation}
    z(x) = h \cos^2\left(\frac{\pi \sqrt{x^2 + y^2}}{2b}\right), \quad \sqrt{x^2 + y^2} \leq b,
\end{equation}
where $h = 0.04~\mathrm{m}$ is the hill height and $b = 0.1~\mathrm{m}$ is the hill half-width. \citet{ishihara1999wind} reported wind-tunnel data at a Reynolds number of $Re = \num{11722}$, based on the hill height and the freestream velocity at the hill crest. To replicate this setup, the present simulation also uses $Re = \num{11722}$, obtained by prescribing a characteristic velocity $U_h = 4.3958~\mathrm{m/s}$ and a kinematic viscosity $\nu = 1.5 \times 10^{-5}~\mathrm{m^2/s}$. The computational domain has dimensions $L_x \times L_y \times L_z = 24h \times 16h \times 8h$. The aerodynamic roughness height in the wind tunnel was reported as $z_0 = 4 \times 10^{-6}~\mathrm{m}$, which is applied uniformly to the wall, with the immersed-boundary surface roughness set to $z_{0,\hbox{\tiny IB}} = 2 z_0$~\cite{fang2016intercomparison, liu2020wall}. The hill is centered at $x = L_x / 4$. Realistic inflow turbulence is generated using a off-line precursor simulation following the approach of \citet{stipa2024tosca}. The precursor domain excludes the hill geometry and employs periodic boundary conditions in the horizontal directions, a wall model at the bottom surface, and a stress-free condition at the top boundary. The resulting flow field is supplied as inflow to the successor simulation, which includes the hill and applies inflow--outflow boundary conditions in the streamwise direction while retaining periodic conditions in the spanwise direction. In the successor domain, the ground surface apart from the hill is represented using an equilibrium wall model, whereas the hill itself is treated with the hybrid wall-modeled immersed boundary method. The stress-free boundary condition is retained at the top boundary.

\subsubsection {Impact of grid resolution}
First, the influence of grid resolution on the streamwise, spanwise, and vertical components of the mean velocity and velocity variance, as well as on the compensated turbulent spectra, are examined. All simulation cases are also compared against the experimental data of \citet{ishihara1999wind}. The mean and variance profiles are extracted along the vertical midplane ($y = 0$) of the computational domain at several streamwise locations ($x/h = -2.5, -1.25, 0.0, 1.25, 2.5, 3.75, 5.0, 6.25$), covering the upstream, hilltop, and lee-side regions, including the separation and reattachment zones. Five different grid resolutions are examined, ranging from coarse to fine: M1 ($96 \times 64 \times 64$) cells, M2 ($120 \times 80 \times 80$) cells, M3 ($144 \times 96 \times 96$) cells, M4 ($192 \times 128 \times 96$) cells, and M5 ($300 \times 200 \times 200$) cells.  

Figure~\ref{fig:avgUWgrid}(a--b) shows the vertical profiles of the mean streamwise and vertical velocity components for the different grid resolutions. In Fig.~\ref{fig:avgUWgrid}(a), the relative difference between the LES and experimental data at $z/h = 1$ varies from $4\%$ to $7\%$, with the finer grids showing slightly larger deviations on the lee side of the hill. For the vertical component, all grids show relative differences below $4\%$ across all locations considered.

\begin{figure}
    \centering
    \begin{tabular}{cc}
    \text{(a)} & \text{(b)} \\
            \includegraphics[width=0.48\linewidth]{plots/3dhill/grid/avgU.pdf} 
            & \includegraphics[width=0.48\linewidth]{plots/3dhill/grid/avgW.pdf} \\
    \end{tabular}
\caption{%
Comparison of different grid resolutions of the hybrid wall-modeled immersed boundary method for flow over a steep three-dimensional hill. The figures show the normalized streamwise and vertical velocity components, $\langle u \rangle /U_\infty + x/h$ and $\langle w \rangle/U_\infty + x/h$, respectively. The vertical axes are scaled by the hill height $h$, and $x/h$ denotes various streamwise locations. The shaded gray region with a black outline represents the hill geometry. Experimental measurements from \citet{ishihara1999wind} are shown as open circles ({\small$\circ$}). Colored solid lines indicate the present LES datasets at different grid resolutions: M1 [\solidline{plotblue}] ($96\times64\times64$), M2 [\dashdotline{plotred}] ($120\times80\times80$), M3 [\dottedline{plotcyan}] ($144\times96\times96$), M4 [\dashedline{plotgreen}] ($192\times128\times96$), and M5 [\solidline{plotgray}] ($300\times200\times200$). A light dashed-gray reference curve [\dashedline{plotgray}] is included to indicate each $x/h$ position.%
}    \label{fig:avgUWgrid}
\end{figure}

Figures~\ref{fig:sigmaUVW-grid}(a--c) compares the vertical profiles of the velocity variance, normalized by the freestream velocity ($U_\infty$), for different grid resolutions. Panels (a), (b), and (c) correspond to the streamwise ($\sigma_u = \sqrt{R_{uu}}/U_{\infty}$), spanwise ($\sigma_v = \sqrt{R_{vv}}/U_{\infty}$), and vertical ($\sigma_w = \sqrt{R_{ww}}/U_{\infty}$) components, respectively. All of the LES grids capture the experimental trends reported by \citet{ishihara1999wind}, showing enhanced turbulence intensity on the lee side of the hill compared with the upwind region. The streamwise normal stress component ($\sigma_u/U_\infty$) reaches its maximum at $z/h = 1$ and $x/h = 1.25$, while the vertical variance ($\sigma_w/U_\infty$) peaks around $z/h = 0.5$. These peaks correspond to strong shear regions produced by flow separation on the lee slope \cite{ishihara1999wind,diebold2013flow,liu2020wall}. Downstream of the reattachment region, all normal stress components gradually decay as the flow recovers toward equilibrium. The maximum values of the normal stress components are listed in Table~\ref{tab:sigma_comparison}. A distinctive feature of the three-dimensional hill wake is the appearance of a secondary local maximum in the spanwise stress ($\sigma_v$) within the near-wall region after $x/h = 2.5$ \cite{ishihara1999wind}, which is absent in the typical two-dimensional wakes reported by \citet{castro1987structure} and \citet{itoh1989turbulence}. In this layer, the $\sigma_u$ profile remains nearly constant, consistent with the observations of \citet{ishihara1999wind} and \citet{castro1982wind}. Among all grids, the streamwise component ($\sigma_u/U_\infty$) peaks on the lee side ($x/h \approx 1.25$--$2.5$), consistent with the separated shear layer. The finer meshes (M4 and M5) slightly overpredict $\sigma_u$ near the crest compared with experimental results, suggesting a mild overestimation of the resolved variance. In contrast, the spanwise and vertical components (Figs.~\ref{fig:sigmaUVW-grid}(b,c)) show weaker dependence on grid resolution, with relative deviations below $10\%$, confirming good convergence for the turbulence variance statistics.

\begin{table*}[t]
\centering
\caption{Comparison of maximum value of normalized stress components ($\sigma_u/U_\infty$, $\sigma_v/U_\infty$, and $\sigma_w/U_\infty$) for different grid resolutions and their relative deviation from the experimental values of \citet{ishihara1999wind}.}
\label{tab:sigma_comparison}
\begin{tabular*}{\textwidth}{@{\extracolsep{\fill}} l l c c c c}
\hline\hline 


Case & Description & $\sigma_u/U_\infty$ & $\sigma_v/U_\infty$ & $\sigma_w/U_\infty$ & Mean error [\%] \\ 
\hline
M1  & 3D hill ($96\times64\times64$)    & 0.291 & 0.141 & 0.208 & 9.7 \\
M2  & 3D hill ($120\times80\times80$)   & 0.291 & 0.150 & 0.200 & 8.9 \\
M3  & 3D hill ($144\times96\times96$)   & 0.338 & 0.173 & 0.210 & 3.6 \\
M4  & 3D hill ($192\times128\times96$)  & 0.328 & 0.173 & 0.207 & 2.7 \\
M5  & 3D hill ($300\times200\times200$) & 0.310 & 0.195 & 0.231 & 4.3 \\
Exp.~\citep{ishihara1999wind} & 3D hill (wind tunnel)       & 0.320 & 0.180 & 0.210 & -- \\
Exp.~\citep{kiya1983structure} & blunt flat plate & 0.250 & 0.170 & 0.220 & -- \\
\hline\hline
\end{tabular*}
\end{table*}

\begin{figure}[htb!]
    \centering
    \text{(a)} \\
    \includegraphics[width=0.52\linewidth]{plots/3dhill/grid/sigmaU_grid.pdf} \\
    \begin{tabular}{cc}
            \text{(b)} & \text{(c)} \\
               \includegraphics[width=0.48\linewidth]{plots/3dhill/grid/sigmaV_grid.pdf} &
                  \includegraphics[width=0.48\linewidth]{plots/3dhill/grid/sigmaW_grid.pdf} \\
    \end{tabular}
\caption{%
Comparison of different grid resolutions of hybrid wall-modeled immersed boundary method for a steep three-dimensional hill. Figures (a), (b), and (c) show the normalized streamwise, spanwise, and vertical velocity variance components, $5\sigma_u/U_\infty + x/h$, $5\sigma_v/U_\infty + x/h$, and $5\sigma_w/U_\infty + x/h$, respectively. Vertical axes are scaled by the hill height $h$, and $x/h$ represents various streamwise locations. The shaded gray region with a black outline indicates the hill geometry. Experimental measurements from \citet{ishihara1999wind} are shown as open circles ({\small$\circ$}). Colored solid lines represent the present LES datasets at different grid resolutions: M1 [\solidline{plotblue}] ($96\times64\times64$), M2 [\dashdotline{plotred}] ($120\times80\times80$), M3 [\dottedline{plotcyan}] ($144\times96\times96$), M4 [\dashedline{plotgreen}] ($192\times128\times96$), and M5 [\solidline{plotgray}] ($300\times200\times200$). A light dashed-gray reference curve [\dashedline{plotgray}] is shown at each $x/h$ position to guide comparison with the reference data.%
}

    \label{fig:sigmaUVW-grid}
\end{figure}
As highlighted by \citet{liu2020wall}, a notable feature of the \citet{ishihara1999wind} dataset is its suitability for evaluating flow statistics on horizontal planes, in addition to the vertical profiles that are more commonly examined in studies of two-dimensional hills or ridges \citep{tian2013experimental,shamsoddin2017large,patton2007turbulent,wan2011large}. Accordingly, the present analysis extends beyond vertical profiles to examine grid-resolution effects on horizontal-plane statistics. Specifically, spanwise profiles of
the mean streamwise velocity and the streamwise normal stress component are analyzed at two representative heights: (i) near the ground ($z/h = 0.125$) and (ii) at $z/h = 1.0$, within
the wake region downstream of the hill at a fixed streamwise location of $x/h = 3.75$. Fig.~\ref{fig:Uspanwise-grid}(a--b) shows the spanwise profiles of the mean streamwise velocity at two vertical heights, $z/h = 0.125$ and $z/h = 1.0$, respectively. Near the ground ($z/h = 0.125$), all LES grids underestimate the streamwise velocity compared to the experimental data. This discrepancy can be attributed to the fact that wall-modeled LES are generally less accurate in the near-surface region, a behavior also observed in three-dimensional hill studies \citep{liu2020wall,diebold2013flow,fang2016intercomparison} and flat-plate LES simulations \citep{stevens2014large}. At $z/h = 1.0$, all cases show excellent agreement with the experimental data, although cases M1, M2, and M4 slightly overpredict the peak location. Similarly, the streamwise velocity variance component is compared at two vertical heights on the horizontal plane at $x/h = 3.75$, as shown in Fig.~\ref{fig:Uspanwise-grid}(c--d). The second-order statistics show excellent agreement across all grid resolutions at both heights, although cases M1, M2, and M4 again slightly overpredict the peak value of the stress component at $z/h = 1.0$.

\begin{figure}[htb!]
    \centering
\begin{tabular}{cc}
  (a) & (b) \\[3pt]
  \includegraphics[width=0.4\linewidth]{plots/3dhill/grid/avgU_spanwisezh125.pdf} &
  \includegraphics[width=0.4\linewidth]{plots/3dhill/grid/avgU_spanwisezh1.pdf} \\
  (c) & (d) \\[3pt]
  \includegraphics[width=0.4\linewidth]{plots/3dhill/grid/sigmaU_spanwisezh125.pdf} &
  \includegraphics[width=0.4\linewidth]{plots/3dhill/grid/sigmaU_spanwisezh1.pdf}  \\
\end{tabular}
\caption{%
Spanwise profiles of the (a,b) streamwise mean velocity $\langle u\rangle/U_\infty$ and (c,d) streamwise velocity variance $\sigma_u/U_\infty$ at a fixed streamwise location $x/h=3.75$ over a steep three-dimensional hill, comparing different grid resolutions of the hybrid wall-modeled immersed boundary method. Panels: (a,c) $z/h=0.125$ and (b,d) $z/h=1.0$. The insets show the hill geometry and mark the sampling position at $x/h=3.75$ and the corresponding height ({\color{red}$\times$}). Experimental measurements from \citet{ishihara1999wind} are shown as open circles ({\small$\circ$}). Colored lines represent the present LES datasets at different grid resolutions: M1 [\solidline{plotblue}] ($96\times64\times64$), M2 [\dottedline{plotcyan}] ($120\times80\times80$), M3 [\dashdotline{plotred}] ($144\times96\times96$), M4 [\dashedline{plotgreen}] ($192\times128\times96$), and M5 [\solidline{plotgray}] ($300\times200\times200$).%
}
\label{fig:Uspanwise-grid}
\end{figure}

% \begin{figure}[htb!]
%     \centering
% \begin{tabular}{cc}
% \textbf{(a)} & \textbf{(b)} \\[3pt]
%       \includegraphics[width=0.3\linewidth]{plots/3dhill/grid/sigmaU_spanwisezh125.pdf} 
%      \includegraphics[width=0.3\linewidth]{plots/3dhill/grid/sigmaU_spanwisezh1.pdf}  \\
% \end{tabular}
% \caption{%
% Spanwise profiles of the streamwise velocity variance $\langle \sigma_u\rangle/U_\infty$ at a fixed streamwise location $x/h=3.75$ over a steep three-dimensional hill, comparing different grid resolutions of the hybrid wall-modeled immersed boundary method. Panels: (a) $z/h=0.125$ and (b) $z/h=1.0$. The insets show the hill geometry and mark the sampling position at $x/h=3.75$ and the corresponding height. Experimental measurements from \citet{ishihara1999wind} are shown as open circles ({\small$\circ$}). Colored lines represent the present LES datasets at different grid resolutions: \textbf{M1} [\solidline{plotblue}] ($96\times64\times64$), \textbf{M2} [\dottedline{plotcyan}] ($120\times80\times80$), \textbf{M3} [\dashdotline{plotred}] ($144\times96\times96$), \textbf{M4} [\dashedline{plotgreen}] ($192\times128\times96$), and \textbf{M5} [\solidline{plotgray}] ($300\times200\times200$). %
% }

% \label{fig:SigmaUspanwise-grid}
% \end{figure}

Finally, we examine the influence of grid resolution using one-dimensional compensated energy spectra. The resolved velocity signals from the LES simulations were recorded at $z/h = 1.0$ in the wake of the hill at $x/h = 3.75$. The sampling frequency was set to $2~\mathrm{kHz}$, which is sufficiently high to capture small-scale motions in the wake of the hill. From Table~\ref{tab:sigma_comparison}, it can be observed that although the turbulence intensity $\sqrt{1/3(\sigma_u^2 + \sigma_v^2 + \sigma_w^2)}/U_\infty$ increases in the lee side of the three-dimensional hill, it remains well below 0.5. Therefore, the \citet{taylor1938spectrum} hypothesis can be applied to convert the frequency spectra to wavenumber space ($k = 2\pi f / U_\infty$) without significant distortion \citep{katul1994intermittency}. The one-dimensional compensated spectra provide a useful means to evaluate the extent of the inertial subrange resolved by each grid. An approximately constant plateau in the compensated spectra is indicative of the presence and extent of the resolved inertial subrange. Fig.~\ref{fig:spectra-grid}(a--c) presents the one-dimensional compensated energy spectra for the streamwise, spanwise, and vertical components for all grid resolutions considered. The results show a clear dependence on grid resolution, with a systematic trend across all components. The finest grid (M5) captures the largest portion of the inertial subrange, as indicated by the broad flat region in the spectra, whereas the coarsest grid (M1) resolves the smallest portion.

\begin{figure}[htb!]
    \centering
   \includegraphics[width=0.8\linewidth]{plots/3dhill/grid/spectraGrid}
    \caption{%
    Comparison of one dimensional compensated energy spectra from large eddy simulations (LES) at different grid resolutions for flow over a steep three-dimensional hill. Panels: (a) streamwise component $E_{uu}(k)$, (b) spanwise component $E_{vv}(k)$, and (c) vertical component $E_{ww}(k)$. The spectra are compensated $k^{5/3}E_{ii}(k)$, where $k$ is the wavenumber. The dashed horizontal line indicates the $-5/3$ Kolmogorov scaling for reference. Colored solid lines represent LES datasets at different grid resolutions: M1 [\solidline{plotblue}] ($96\times64\times64$), M2 [\dottedline{plotcyan}] ($120\times80\times80$), M3 [\dashdotline{plotred}] ($144\times96\times96$), M4 [\dashedline{plotgreen}] ($192\times128\times96$), and M5 [\solidline{plotgray}] ($300\times200\times200$).%
    }
    \label{fig:spectra-grid}
\end{figure}

Overall, the grid-resolution study demonstrates that all cases reproduce the main experimental features of the flow, including the velocity speed-up over the hill, the location of the recirculation zone, and the distribution of turbulence intensity. While the finest grid (M5) captures a broader inertial subrange, the improvement over the intermediate grids (M3 and M4) is marginal compared to the substantial increase in computational cost. The M3 grid ($144\times96\times96$) provides an excellent compromise between accuracy and efficiency. It accurately reproduces both the mean-flow characteristics and second-order statistics with deviations below $5\%$ from the experimental data, while remaining computationally affordable for extended parametric studies. Therefore, the M3 resolution is adopted for subsequent analyses in this work.

\subsubsection {Comparison with other wall-modeled immersed boundary methods} 

The hybrid wall-modeled immersed boundary method used in the present study is closely aligned with the Chester wall-modeled immersed boundary approach~\cite{chester2007modeling,liu2020wall}, with the key distinction being an additional velocity reconstruction procedure for the ghost cells. Therefore, it is of particular interest to compare the present method with similar approaches. All the datasets considered correspond to large eddy simulations of the \citet{ishihara1999wind} wind-tunnel experiment. Table~\ref{tab:datasetsLES} summarizes these datasets, which employ different wall-modeled IBM formulations and subgrid-scale models. Although the LES domains vary slightly among the datasets, grid-resolution metrics are provided to facilitate a consistent comparison. From Table~\ref{tab:datasetsLES}, it is evident that the present M3 case uses one of the coarsest grids, being approximately four times coarser than the finest resolution of \citet{liu2020wall} and twice as coarse as their coarse-resolution case.

\begin{table}[htb!]
\centering
\caption{Comparison of grid resolutions used in wall-modeled immersed boundary method datasets. Here, LASD denotes the Lagrangian scale-dependent subgrid-scale model~\cite{bou2005scale}, and NDC refers to the normal-derivative correction  for the velocity gradient. The Smagorinsky model employed in \citet{liu2020wall} includes a phenomenological wall-damping function proposed by \citet{mason1992stochastic}.}
\label{tab:resolution_comparison}
\resizebox{\textwidth}{!}{%
\begin{tabular}{l l l l S[table-format=1.2] S[table-format=1.2] S[table-format=1.2]}
\hline\hline
Dataset & IBM Method & Numerical Code & SGS Model &
{$\Delta x~(\mathrm{mm})$} & {$\Delta y~(\mathrm{mm})$} & {$\Delta z~(\mathrm{mm})$} \\
\hline
\citet{liu2020wall} (Coarse) & Simplified Chester model~\cite{chester2007modeling,liu2020wall} & Pseudo-spectral & Smagorinsky & 3.30 & 3.30 & 1.65 \\
\citet{liu2020wall} (Fine)   & Simplified Chester model~\cite{chester2007modeling,liu2020wall} & Pseudo-spectral & Smagorinsky & 1.65 & 1.65 & 0.83 \\
\citet{diebold2013flow}      & Smearing method~\cite{diebold2013flow} & Pseudo-spectral (EPFL-LES) & LASD & 2.50 & 2.50 & 2.50 \\
\citet{fang2016intercomparison} & Smearing method with NDC~\cite{fang2016intercomparison} & Pseudo-spectral & LASD & 2.50 & 2.50 & 1.25 \\
\citet{ishihara1999wind} & Wind-tunnel & -- & -- & {\text{--}} & {\text{--}} & {\text{--}} \\

Present (M3) & Hybrid wall-modeled (sharp-interface) & Finite-volume (TOSCA) & BAMD & 6.67 & 6.67 & 3.33 \\
\hline\hline
\end{tabular}%
}
\label{tab:datasetsLES}
\end{table}
%
Fig.~\ref{fig:lesdataU} compares the vertical profiles of the mean streamwise velocity obtained from the different LES datasets listed in Table~\ref{tab:datasetsLES}. All models show negligible differences in predicting the mean flow behavior around the hill, including the upwind, hilltop, and lee-side regions. However, near the ground on the lee side, all models tend to underpredict the mean streamwise velocity compared with the experimental dataset. This discrepancy can be attributed primarily to the limitations of wall-modeled approaches in regions of strong flow separation and recirculation \cite{liu2020wall,stevens2014large}. Among all datasets, \citet{diebold2013flow} exhibits the largest deviation in the mean streamwise velocity. Notably, the method presented in this study is well aligned with all previous LES results, although with a much coarser grid resolution. Fig.~\ref{fig:lesdataSigmaU}(a--b) shows the streamwise normal stress component from various LES datasets. As seen in Fig.~\ref{fig:lesdataSigmaU}, the LES by \citet{liu2020wall} tends to overestimate the streamwise normal stress, both upstream and at the hill crest. Their results also display pronounced near-wall oscillations. To mitigate these spurious fluctuations and achieve closer agreement with experimental measurements, a substantially finer grid resolution was adopted together with pseudo-spectral discretization in the horizontal directions. The oscillations primarily appear around the hilltop, where strong horizontal gradients develop, and although the higher resolution reduced their amplitude, they did not disappear entirely. Interestingly, in the lee-side separation zone, all fine-resolution datasets show a distinct jump in velocity variance, as highlighted in the zoomed-in view in Fig.~\ref{fig:lesdataSigmaU}(b). \citet{fang2016intercomparison} and \citet{diebold2013flow} report a much larger jump compared with the present fine-grid case (M5) and \citet{liu2020wall}. This difference is likely related to the fact that flow separation is inherently a viscous phenomenon, while wall-modeled approaches do not resolve the viscous sublayer. Increasing the vertical resolution can partially capture this layer, but it may still yield inaccurate shear-stress predictions from the wall model. However, during the grid-resolution study, this jump in velocity variance was not observed in cases M1, M2, and M3 (Fig.~\ref{fig:sigmaUVW-grid}(a--c)); the feature appears and becomes more pronounced only as the grid is refined.
%
\begin{figure}[htb!]
    \centering
\begin{tabular}{c}
      \includegraphics[width=0.55\linewidth]{plots/3dhill/explesdatas/avgULESdata.pdf} 
\end{tabular}
\caption{Comparison of wall-modeled immersed boundary methods for atmospheric boundary layer (ABL) flow over a steep three-dimensional hill. The plot shows the streamwise mean velocity $\langle u\rangle/U_\infty + x/h$, and vertical axes are scaled by the hill height $h$ and $x/h$ represents various streamwise locations. The shaded gray region with a black outline indicates the hill geometry. Experimental measurements from \citet{ishihara1999wind} are plotted as open circles $(\circ)$. Colored solid lines denote LES datasets using different wall-modeled immersed boundary formulations: present method M3 (\colorline{plotblue}), \citet{diebold2013flow} (\colorline{plotred}), \citet{liu2020wall} (\colorline{plotcyan}), \citet{fang2016intercomparison} (\colorline{plotgreen}), and a light dashed-gray reference curve (\dashedline{plotgray}) is shown at each $x/h$ position to guide comparison with reference data.}
\label{fig:lesdataU}
\end{figure}

\begin{figure}[htb!]
  \centering
  \begin{tabular}{@{}c@{\hspace{0.8em}}c@{}}
  $(a)$ & $(b)$ \\
    \includegraphics[width=0.55\linewidth]{plots/3dhill/explesdatas/sigmaU_LESdata_full.pdf} &
    \raisebox{2.5mm}{\includegraphics[width=0.35\linewidth]{plots/3dhill/explesdatas/sigmaU_LESdata_zoomed.pdf}}
  \end{tabular}
\caption{Comparison of wall-modeled immersed boundary methods for atmospheric boundary layer (ABL) flow over a steep three-dimensional hill. 
  Panel (a) shows the complete distribution of streamwise stress component $5\sigma_u/U_\infty + x/h$, and panel (b) provides a zoomed-in view of the shaded gray rectangular region in (a). Vertical axes are normalized by hill height $h$, and $x/h$ represents different streamwise positions. The shaded gray area with a black outline denotes the hill geometry. Experimental data from \citet{ishihara1999wind} are shown as open circles $(\circ)$, while colored lines correspond to LES datasets using different wall-modeled immersed boundary formulations: present method \textbf{M3} (\colorline{plotblue}), \citet{diebold2013flow} (\colorline{plotred}), \citet{liu2020wall} (\colorline{plotcyan}), \citet{fang2016intercomparison} (\colorline{plotgreen}), and a dashed gray reference curve (\dashedline{plotgray}) for visual comparison.}
\label{fig:lesdataSigmaU}
\end{figure}
%
Figure~\ref{fig:spanwiseU-lesdata}(a--b) presents the spanwise profiles of the mean streamwise velocity $\langle u \rangle /U_\infty$ at two vertical locations above the hill surface. Panel~(a) corresponds to $z/h = 0.125$, close to the near-wall region, while panel~(b) represents $z/h = 1.0$ at top of hill. In both cases, the streamwise location is fixed at $x/h = 3.75$, corresponding to the downstream side of the hill. All LES datasets consistently underpredict the streamwise velocity near the surface compared to the wind-tunnel data of \citet{ishihara1999wind}, a behavior commonly linked to the wall-model formulation. This trend is evident across all datasets considered. Among the cases reported by \citet{liu2020wall}, the coarse-grid simulation shows the strongest velocity deficit, while their fine-grid case exhibits close agreement with the present LES (M3). As shown in Fig.~\ref{fig:spanwiseU-lesdata}a, both the present M3 case and the fine-grid LES of \citet{liu2020wall} reproduce the near-wall velocity distribution reasonably well. At the hill height ($z/h = 1.0$), shown in Fig.~\ref{fig:spanwiseU-lesdata}(b), all datasets converge toward the experimental profile, except for \citet{diebold2013flow}, which still underpredicts the velocity relative to other simulations.

Similarly, we compare the streamwise stress component on horizontal planes across various LES and wind-tunnel datasets, as shown in Fig.~\ref{fig:spanwiseU-lesdata}(c--d). At $z/h = 0.125$, the differences among the LES datasets are negligible. However, minor deviations appear near the hilltop, where the coarse-grid case of \citet{liu2020wall} slightly underpredicts the velocity deficit, while the finest-grid case \cite{liu2020wall} and the present M3 simulation aligns closely with the wind tunnel data \cite{ishihara1999wind}. 
%
\begin{figure}
    \centering
\begin{tabular}{cc}
\text{(a)} & \text{(b)} \\
      \includegraphics[width=0.4\linewidth]{plots/3dhill/explesdatas/avgU_spanwise_lesdata_zh125.pdf} &
     \includegraphics[width=0.4\linewidth]{plots/3dhill/explesdatas/avgU_spanwise_lesdata_zh1.pdf}  \\
     \text{(c)} & \text{(d)} \\
      \includegraphics[width=0.4\linewidth]{plots/3dhill/explesdatas/sigmaU_spanwise_lesdata_zh125.pdf}& 
     \includegraphics[width=0.4\linewidth]{plots/3dhill/explesdatas/sigmaU_spanwise_lesdata_zh1.pdf}  \\
\end{tabular}
\caption{%
Spanwise profiles of the (a,b) streamwise velocity $\langle u \rangle/U_\infty$ and (c,d) streamwise velocity variance $\sigma_U/U_\infty$ at a fixed streamwise location $x/h=3.75$ over a steep three-dimensional hill, comparing different wall-modeled immersed boundary methods. Panels: (a,c) $z/h=0.125$ and (b,d) $z/h=1.0$. The insets show the hill geometry (gray shaded with black outline) and mark the sampling position at $x/h=3.75$ and the corresponding height ({\color{red}$\times$}). Experimental measurements from \citet{ishihara1999wind} are shown as open circles ({\small$\circ$}). Colored solid lines represent LES datasets obtained using various wall-modeled immersed boundary formulations: present method M3 (\colorline{plotblue}), coarse mesh ($384 \times 192 \times 192$) from \citet{liu2020wall} (\colorline{plotred}), fine mesh ($768 \times 384 \times 385$) from \citet{liu2020wall} (\colorline{plotcyan}),\citet{fang2016intercomparison} (\colorline{plotgreen}), and \citet{diebold2013flow} (\colorline{plotgray}).%
 }
\label{fig:spanwiseU-lesdata}
\end{figure}

% \begin{figure}
%     \centering
% \begin{tabular}{cc}
% \text{(a)} & \text{(b)} \\
%       \includegraphics[width=0.3\linewidth]{plots/3dhill/explesdatas/sigmaU_spanwise_lesdata_zh125.pdf}& 
%      \includegraphics[width=0.3\linewidth]{plots/3dhill/explesdatas/sigmaU_spanwise_lesdata_zh1.pdf}  \\
% \end{tabular}
% \caption{%
% Spanwise profiles of the streamwise velocity variance $\langle \sigma_U \rangle/U_\infty$ at a fixed streamwise location $x/h=3.75$ over a steep three-dimensional hill, comparing different wall-modeled immersed boundary methods. Panels: (a) $z/h=0.125$ and (b) $z/h=1.0$. The insets show the hill geometry (gray shaded with black outline) and mark the sampling position at $x/h=3.75$ and the corresponding height. Experimental measurements from \citet{ishihara1999wind} are shown as open circles ({\small$\circ$}). Colored solid lines represent LES datasets obtained using various wall-modeled immersed boundary formulations: present method \textbf{M3} (\colorline{plotblue}), coarse mesh ($384 \times 192 \times 192$) from \citet{liu2020wall} (\colorline{plotred}), fine mesh ($768 \times 384 \times 385$) from \citet{liu2020wall} (\colorline{plotcyan}), and \citet{fang2016intercomparison} (\colorline{plotgreen}).%
% }
% \label{fig:spanwiseSigmaU-lesdata}
% \end{figure}

\subsubsection {Characterization of flow structures}

In this section, we assess how effectively the hybrid wall-modeled immersed boundary method reproduces the dominant flow structures in atmospheric flow over a three-dimensional hill. Although the $\mathcal{Q}$- and $\lambda_2$-criteria~\cite{jeong1995identification} are commonly used to visualize coherent structures in such configurations, at high Reynolds numbers these techniques reveal a multitude of small-scale features, making it difficult to distinguish the large-scale coherent motions even after thresholding. The choice of the threshold value plays a crucial role in isolating relevant vortical structures: a low threshold tends to include a large number of weak, small-scale eddies, which obscures the organized vortical patterns, while a high threshold can remove physically significant structures by retaining only the strongest cores. Consequently, even with careful tuning, the $\mathcal{Q}$- and $\lambda_2$-criteria may not clearly represent the dominant coherent features in high-Reynolds-number flows.  

An alternative approach is the proper orthogonal decomposition (POD), but as noted by~\citet{garcia2009large}, it is computationally prohibitive to obtain converged fields even for the first few modes. Following~\citet{garcia2006identification,frohlich2005highly,garcia2009large}, we adopt their method of using streamwise velocity fluctuations to identify streaks, noting that coherent structures are also known to coincide with pressure minima~\cite{jeong1995identification,dubief2000coherent}. Pressure fluctuation fields have been shown to reveal coherent vortical regions more clearly than instantaneous pressure because the latter is influenced by the spatially varying mean pressure field that is unrelated to instantaneous motions; this effect is eliminated by subtracting the mean pressure, $\langle p \rangle$.  

Figure~\ref{fig:3dhill_hairpin}(a--b) shows iso-surfaces of fluctuating velocity and pressure fields illustrating coherent structure formation in the lee of the hill. Red and blue surfaces represent high- and low-speed streaks ($u' = +0.35~\mathrm{m/s}$ and $u' = -0.35~\mathrm{m/s}$), while the green iso-surface ($p' = -0.1~\mathrm{m^2/s^2}$) marks regions of low pressure associated with vortex cores. Our focus is on identifying hairpin vortices rather than individual streaks; therefore, a relatively high threshold for the streamwise velocity fluctuation is used. The three-dimensional view (top) reveals the characteristic head--leg pattern of hairpin vortices forming downstream of the crest, while the side view (bottom) highlights their spatial organization and vertical extent within the near-surface region.  

\begin{figure}[htb!]
    \centering
    \begin{tabular}{c}
        \includegraphics[width=0.8\linewidth]{plots/3dhill/hillQ2-edited.pdf} \\
        \vspace{4em}
        \includegraphics[width=0.8\linewidth]{plots/3dhill/hillQ2-xz.pdf}
    \end{tabular}
    \caption{Iso-surfaces of the fluctuating velocity and pressure fields illustrating the formation of coherent structures in the lee of the hill. 
    Red and blue surfaces correspond to positive and negative streamwise velocity fluctuations ($u' = +0.35~\mathrm{m/s}$ and $u' = -0.35~\mathrm{m/s}$), while the green iso-surface ($p' = -0.1~\mathrm{m^2/s^2}$) highlights regions of low pressure associated with vortex cores. The three-dimensional view (top) shows the characteristic head--leg pattern of hairpin vortices forming downstream of the hill crest, whereas the side view (bottom) emphasizes their spatial organization and vertical extent within the near-surface region.}
    \label{fig:3dhill_hairpin}
\end{figure}

Having established the qualitative organization of coherent structures, we now examine their statistical signatures through the turbulent energy spectra. Previous large eddy simulations of flow over three-dimensional hills~\cite{diebold2013flow,fang2016intercomparison,liu2020wall} have not investigated spectral characteristics, whereas the wind-tunnel experiments of~\citet{ishihara1999wind} provided detailed spectral measurements. It is therefore interesting to evaluate how the turbulence structure in the hill wake is represented when the terrain is modeled using an immersed boundary approach.  Figure~\ref{fig:explesspectra}(a--c) compares the premultiplied, variance-normalized turbulent energy spectra of the streamwise, spanwise, and vertical velocity components between the present LES (M3) and the wind-tunnel data, using the Strouhal scaling $(f h / U_h)$, where $f$ is the frequency and $h$ is the hill height. The spectra are estimated using the maximum entropy method (MEM)~\cite{burg1975maximum}, following the same approach used by~\citet{ishihara1999wind}.  

The streamwise spectrum at $z/h = 1$ [Fig.~\ref{fig:explesspectra}(a)] shows a broad peak near $f h / U_h \approx 0.13$, in agreement with the reference data, corresponding to the vortex shedding frequency from the separation bubble. The spanwise spectrum at $z/h = 0.125$ in Fig.~\ref{fig:explesspectra}(b) exhibits a peak at about half that frequency, consistent with low-frequency motions in the wall layer behind the hill. The vertical spectrum at $z/h = 0.5$ in Fig.~\ref{fig:explesspectra}(c) also displays a peak near $f h / U_h \approx 0.13$, indicating that the vortices shed from the separation bubble are inclined downward due to the descending flow in the lee of the hill. These results indicate that the LES with the immersed boundary method reproduces the dominant spectral features and peak locations reported in the wind-tunnel measurements, demonstrating that the approach captures the primary unsteady flow dynamics in the hill wake.

\begin{figure}[htb!]
    \centering
 
     \includegraphics[width=0.95\linewidth]{plots/3dhill/explesdatas/MEM_Spectra_Final}


\caption{%
Comparison of one-dimensional premultiplied energy spectra from the present LES at various vertical positions $z/h$ for the hybrid wall-modeled immersed boundary method applied to flow over a steep three-dimensional hill. Panels show (a) the streamwise component $E_{uu}(f)$, (b) the spanwise component $E_{vv}(f)$, and (c) the vertical component $E_{ww}(f)$. The spectra are premultiplied by frequency $f$ and normalized by the corresponding velocity variance $\sigma_i^2$ as $fE_{ii}/\sigma_i^2$. The horizontal axis represents the non-dimensional frequency $n h / U_h$, where $h$ is the hill height and $U_h$ is the mean velocity at hill height, characterizing the ratio of turbulent to advective time scales. Experimental data from \citet{ishihara1999wind} are included for comparison: the streamwise spectra at $z/h=1.0$, the spanwise spectra at $z/h=0.125$, and the vertical spectra at $z/h=0.5$, all computed at the probe location $x/h=3.75$.}

\label{fig:explesspectra}
\end{figure}

\subsubsection{Impact of turbulence modelling}
In this section, we evaluate the influence of subgrid-scale (SGS) modelling on the prediction of key flow statistics and small-scale turbulence characteristics. The performance of recent dynamic models, including the LASD, LASI, and localized DSM, is compared against the mixed subgrid-scale formulation based on the BAMD approach. The comparison focuses on several aspects: mean velocity distributions, normal stress components, and spanwise profiles, followed by a spectral analysis of the one-dimensional compensated turbulent energy spectra. 

First, we compared the mixed modeling approach with the dynamic class of eddy-viscosity models based on the vertical profiles of mean and normal stress components. Figs.~\ref{fig:sigmaUW_les}(a--b) present the vertical distributions of the mean streamwise velocity predicted by different SGS models. At the coarse resolution (Fig.~\ref{fig:sigmaUW_les}a), the dynamic eddy-viscosity models show differences from the experimental reference on the lee side of the hill, where turbulence is affected by flow separation, recirculation, and reattachment. These regions modify the local shear and lead to noticeable variations in the mean velocity profile from the surface up to the hill height. Among the dynamic models, the LASD formulation exhibits the largest variation, which may be related to its two-level test filtering and the associated near-wall resolution requirements. The DSM and LASI models display a similar trend in these regions. This behavior is consistent with the general observation that most existing SGS models are formulated under the assumption of locally isotropic subgrid motions and negligible numerical or grid-induced inhomogeneity~\cite{moser2021statistical,rozema2015minimum}. In separated or strongly anisotropic flows, these assumptions become less valid, and the flow-model interaction becomes sensitive to grid resolution. As shown in Fig.~\ref{fig:sigmaUW_les}(b), finer grids improve the representation of the mean profile for the dynamic eddy-viscosity models, suggesting that part of the coarse-grid deviation is attributable to limited grid resolution rather than the SGS formulation itself.
\begin{figure}
    \centering
    \begin{tabular}{cc}
        \text{(a) \footnotesize Coarse Resolution} & \text{(b) \footnotesize Fine Resolution} \\
            \includegraphics[width=0.4\linewidth]{plots/3dhill/lescomparison/avgU_LES_coarse.pdf} 
            & \includegraphics[width=0.4\linewidth]{plots/3dhill/lescomparison/avgU_LES_fine.pdf} \\
    \text{(c) \footnotesize Coarse resolution} & \text{(d) \footnotesize Fine resolution} \\
            \includegraphics[width=0.4\linewidth]{plots/3dhill/lescomparison/sigmaU_lesmodel_coarse.pdf} 
            & \includegraphics[width=0.4\linewidth]{plots/3dhill/lescomparison/sigmaU_lesmodel_fine.pdf} \\
        \text{(e) \footnotesize Coarse resolution} & \text{(f) \footnotesize Fine resolution} \\
            \includegraphics[width=0.4\linewidth]{plots/3dhill/lescomparison/sigmaW_lesmodel_coarse.pdf} 
            & \includegraphics[width=0.4\linewidth]{plots/3dhill/lescomparison/sigmaW_lesmodel_fine.pdf} \\
    \end{tabular}
\caption{%
Comparison of SGS models at coarse ($144\times96\times96$) and fine ($300\times200\times200$) grid resolutions for flow over a three-dimensional hill using the hybrid wall-modeled immersed boundary method. Panels: (a,b) streamwise velocity component $\langle u \rangle/U_\infty$, (c,d) streamwise velocity variance $5\sigma_u/U_\infty + x/h$ and (e,f) vertical velocity variance $5\sigma_w/U_\infty + x/h$. In the fine-resolution case, the BAMD model is retained at the coarse resolution for reference. Experimental measurements from~\citet{ishihara1999wind} are shown as open circles~({\small$\circ$}). Colored lines represent different SGS models: BAMD [\solidline{plotblue}], DSM [\dashdotline{plotcyan}], LASI [\dottedline{plotred}], and LASD [\dashedline{plotgreen}].%
}
\label{fig:sigmaUW_les}
\end{figure}

The BAMD mixed model combines a scale-similarity structural term with a dissipative term based on the AMD concept, in which the eddy viscosity is computed from a fluctuating scalar eddy viscosity derived from the triple product of velocity gradients. This approach incorporates grid-induced effects and this formulation naturally limits excessive dissipation and avoids numerical instability in regions with strong velocity gradients, such as shear layers or separation zones. However, as noted by \citet{kamal2024artificial}, the AMD formulation excludes the explicit contribution of vortex stretching, which is considered as the primary mechanism for forward energy cascade. Minor differences observed near the wall on the lee side may reflect both this simplification and the limitations of the wall-model in separated flows. When the grid is refined to roughly twice the coarse resolution (Fig.~\ref{fig:sigmaUW_les}b), all dynamic models exhibit closer agreement and begin to converge, indicating improved resolution of the near-surface dynamics. At sufficiently fine grids, the sensitivity to the SGS model becomes weak, consistent with previous studies~\citep[e.g.,][]{meneveau2000scale}. Similar behavior was observed in the vertical profiles of the vertical velocity component, which are not shown here for brevity.

Figs.~\ref{fig:sigmaUW_les}(c--f) compare the streamwise and vertical velocity variances across the different SGS models and grid resolutions. All models capture the overall distribution of turbulence intensity, showing close agreement with the measurements upwind and near the hill crest, where the flow remains attached. The strongest variations appear in the wake region, reflecting the impact of separation and reattachment on turbulence anisotropy. At the coarse resolution, the mixed model tends to reproduce the peak stress levels more accurately, while the dynamic eddy-viscosity models show slightly lower variance in the separated shear layer. With grid refinement, the predicted stress magnitudes increase and align closely with the experimental profiles, indicating improved resolution of the shear layer and recirculation zone. Similar trends are observed for the vertical component, confirming that model differences diminish as the grid becomes sufficiently fine.

Figure~\ref{fig:avgUspanwiselesmodel}(a--d) show spanwise profiles of the normalized mean streamwise velocity at $x/h = 3.75$ over the steep three-dimensional hill. Results are presented at two vertical locations, $z/h = 0.125$ and $z/h = 1.0$, using coarse (M3) and fine (M5) grid resolutions. Experimental measurements from \citet{ishihara1999wind} are included for reference. At the near-surface location $z/h = 0.125$, all models exhibit a clear wake deficit at the coarse grid resolution (Fig.~\ref{fig:avgUspanwiselesmodel}a). DSM, LASI and LASD predict a higher wake deficit and a smeared peak at the centerline compared to the experimental data over the central region of the profile. The BAMD model shows closer agreement with the measurements in both the peak magnitude and the
lateral extent of the wake deficit. The LASD results show an additional deviation near the hill centerline. With grid refinement (Fig.~\ref{fig:avgUspanwiselesmodel}b), the mean velocity profiles converge toward the experimental data and the agreement among models improves, although all models continue to predict a slightly higher wake deficit. The BAMD results, shown at the coarse grid resolution in both panels, remain comparable to the fine-resolution results obtained with the other models. At $z/h = 1.0$, the coarse-grid results (Fig.~\ref{fig:avgUspanwiselesmodel}c) indicate that DSM, LASI and LASD predict a reduced wake deficit relative to the experimental measurements, whereas the BAMD results exhibit excellent agreement with the measurements across the span. At the fine resolution (Fig.~\ref{fig:avgUspanwiselesmodel}d), the profiles from DSM and LASI converge toward the experimental data, while a centerline deviation persists in the LASD results.

\begin{figure}[htb!]
    \centering
    \begin{tabular}{cc}
    \text{(a) \footnotesize Coarse resolution} & \text{(b) \footnotesize Fine resolution} \\
            \includegraphics[width=0.4\linewidth]{plots/3dhill/lescomparison/avgU_spanwisezh25_coarse_LESModels.pdf} &  
           \includegraphics[width=0.4\linewidth]{plots/3dhill/lescomparison/avgU_spanwisezh25_fine_LESModels.pdf}  \\
    \text{(c) \footnotesize Coarse resolution} & \text{(d) \footnotesize Fine resolution} \\
            \includegraphics[width=0.4\linewidth]{plots/3dhill/lescomparison/avgU_spanwisezh1_coarse_LESModels.pdf} &  
           \includegraphics[width=0.4\linewidth]{plots/3dhill/lescomparison/avgU_spanwisezh1_fine_LESModels.pdf}  \\
           
    \end{tabular}

\caption{%
Comparison of spanwise profiles of the normalized streamwise velocity $\langle u \rangle/U_\infty$ obtained using different subgrid-scale models at a fixed streamwise location, $x/h=3.75$, over a steep three-dimensional hill. Results are shown at two vertical locations: $z/h=0.125$ and $z/h=1.0$. The insets illustrate the hill geometry (gray shaded with a black outline) and indicate the sampling location at $x/h=3.75$ and the corresponding height ({\color{red}$\times$}). Panel (a) shows the comparison at $z/h=0.125$ for all models using the coarse grid resolution (M3, $144\times96\times96$). Panel (b) shows the same comparison using a finer grid resolution (M5, $300\times200\times200$), with the BAMD model at the coarse resolution ($144\times96\times96$) retained for reference. Panels (c) and (d) show the corresponding comparisons at $z/h=1.0$. Experimental measurements from~\citet{ishihara1999wind} are shown as open circles~({\small$\circ$}). Colored lines represent different SGS models: BAMD [\solidline{plotblue}], DSM [\dashdotline{plotcyan}], LASI [\dottedline{plotred}], and LASD [\dashedline{plotgreen}].%
}
    \label{fig:avgUspanwiselesmodel}
\end{figure}

To assess how these differences in the mean flow are reflected in second-order statistics, the corresponding spanwise profiles of the streamwise velocity variance are examined in Fig.~\ref{fig:sigmaUspanwiselesmodel}(a--d). At the near-surface location $z/h = 0.125$, the coarse-grid results show clear differences among the SGS models in both the magnitude and spanwise distribution of the streamwise velocity variance. DSM and LASI predict broader variance peaks near the centerline compared to the experimental data, while LASD exhibits more strongly smeared variance peaks with a flatter spanwise distribution. The BAMD model provides closer agreement with the measurements, capturing both the peak level and the lateral extent of the variance distribution. With grid refinement, the predicted variance profiles converge toward the experimental data and the spread among models is reduced. At $z/h = 1.0$, the coarse-grid results show larger discrepancies among the SGS models, particularly in the wake region. DSM and LASI underpredict the peak level of the streamwise velocity variance, while LASD exhibits a laterally shifted distribution relative to the measurements. At the fine resolution, the profiles from DSM and LASI converge toward the experimental data.

\begin{figure}[htb!]
    \centering
    \begin{tabular}{cc}
    \text{(a) Coarse Resolution} & \text{(b) Fine Resolution} \\
            \includegraphics[width=0.4\linewidth]{plots/3dhill/lescomparison/sigmaU_coarse_spanwisezh125.pdf} &  
           \includegraphics[width=0.4\linewidth]{plots/3dhill/lescomparison/sigmaU_fine_spanwisezh125.pdf}  \\
\text{(c) Coarse Resolution} & \text{(d) Fine Resolution} \\
            \includegraphics[width=0.4\linewidth]{plots/3dhill/lescomparison/sigmaU_coarse_spanwisezh1.pdf} &  
           \includegraphics[width=0.4\linewidth]{plots/3dhill/lescomparison/sigmaU_fine_spanwisezh1.pdf}  \\
           
    \end{tabular}

\caption{%
Comparison of spanwise profiles of the normalized streamwise velocity variance $\sigma_U/U_\infty$ obtained using different subgrid-scale models at a fixed streamwise location, $x/h=3.75$, over a steep three-dimensional hill. Results are shown at two vertical locations: $z/h=0.125$ and $z/h=1.0$. The insets illustrate the hill geometry (gray shaded with a black outline) and indicate the sampling location at $x/h=3.75$ and the corresponding height ({\color{red}$\times$}). Panel (a) shows the comparison at $z/h=0.125$ for all models using the coarse grid resolution (M3, $144\times96\times96$). Panel (b) shows the same comparison using a finer grid resolution (M5, $300\times200\times200$), with the BAMD model at the coarse resolution ($144\times96\times96$) retained for reference. Panels (c) and (d) show the corresponding comparisons at $z/h=1.0$. Experimental measurements from~\citet{ishihara1999wind} are shown as open circles~({\small$\circ$}). Colored lines represent different SGS models: BAMD [\solidline{plotblue}], DSM [\dashdotline{plotcyan}], LASI [\dottedline{plotred}], and LASD [\dashedline{plotgreen}].%
}
    \label{fig:sigmaUspanwiselesmodel}
\end{figure}

\begin{figure}[htb!]
    \centering
    \begin{tabular}{cc}
    \text{(a) \footnotesize Coarse Resolution} & \text{(b) \footnotesize Fine Resolution} \\
            \includegraphics[height=0.35\linewidth]{plots/3dhill/lescomparison/gamma_coarse_LES.pdf} &  
           \includegraphics[height=0.35\linewidth]{plots/3dhill/lescomparison/gamma_fine_LES.pdf}  \\
           
    \end{tabular}
\caption{Comparison of the fraction of total turbulent kinetic energy, $\gamma$, across different subgrid-scale models. Panel (a) shows results at coarse grid resolution for the M3 case ($144\times96\times96$). Panel (b) shows results for the M5 case at fine grid resolution, while the BAMD model is retained at coarse resolution (M3) for reference.}
\label{fig:gammalesmodel}
\end{figure}

Finally, we compute the vertical profiles of the fraction of resolved turbulent kinetic energy, $\gamma(z)$, which is defined as
\begin{equation}
\gamma(z) \;=\; \frac{k_{\mathrm{res}}(z)}{k_{\mathrm{res}}(z) + k_{\mathrm{sgs}}(z)},
\end{equation}
where $k_{\mathrm{res}}(z) = \tfrac{1}{2}\langle \overline {u_i' u_i'} \rangle$ denotes the resolved turbulent kinetic energy and $k_{\mathrm{sgs}}(z)$ is the modeled subgrid-scale turbulent kinetic energy. This quantity represents the fraction of the total turbulent kinetic energy resolved by the LES at a given height. Following \citet{Pope_2000}, the resolved energy fraction is used as a diagnostic indicator of LES resolution quality rather than as a strict criterion. Larger values of $\gamma$ indicate that a greater fraction of the total turbulent kinetic energy is resolved and that the subgrid-scale model contributes primarily at the smallest unresolved scales. No universal threshold exists \cite{wurps2020grid, davidson2009large, singh2023large}; in particular, even a coarse LES can yield $\gamma \gtrsim 0.9$ depending on the flow and the relative contributions of resolved and modeled kinetic energy. Therefore, the vertical dashed gray line at $\gamma = 0.95$ in Fig.~\ref{fig:gammalesmodel} is shown only as a reference level
to indicate cases in which nearly all of the turbulent kinetic energy is resolved, and it is used here for comparative purposes across different subgrid-scale models.

At the coarse grid resolution (Fig. \ref{fig:gammalesmodel}a), DSM, LASI, and LASD predict lower values of $\gamma(z)$, with LASD exhibiting the largest deviation and a more non-uniform vertical distribution throughout the boundary layer. In contrast, the BAMD model predicts values of $\gamma(z)$ close to unity over most of the vertical extent. At the fine grid resolution (Fig. \ref{fig:gammalesmodel}b), the profiles from DSM, LASI, and LASD shift toward higher values of $\gamma(z)$. The BAMD results, computed at the coarse grid resolution, remain comparable to the fine-grid predictions from the other models. %Overall, grid refinement improves the prediction of streamwise velocity and its variance profiles for LASD, LASI, and DSM models, while the BAMD model exhibit remarkable agreement with experimental wind tunnel data at relatively coarse grid resolutions.

\subsection{Realistic Terrain: Bolund Hill}
To assess the capability of the proposed hybrid wall-modeled immersed boundary framework combined with a mixed SGS model under realistic atmospheric conditions, we consider an ABL flow over the Bolund hill, a well-documented natural terrain located in Roskilde Fjord, Denmark. Bolund is an isolated, grass-covered hill approximately $12\,\mathrm{m}$ high, with characteristic horizontal dimensions of about $130\,\mathrm{m}$ in length and $75\,\mathrm{m}$ in width. The site is surrounded by water on three sides, while a narrow, intermittently submerged isthmus connects it to the mainland on the eastern side as shown in Fig. \ref{fig:realbolund}. This configuration creates an abrupt transition from low-roughness water to higher-roughness land, making Bolund a challenging test case for numerical simulations of atmospheric boundary layer flow over complex terrain \cite{bechmann2011bolund,berg2011bolund,diebold2013flow,prospathopoulos2012application,conan2016experimental,schubiger2020evaluation}.

\begin{figure}[htb!]
	\includegraphics[width=0.6\linewidth]{plots/bolundhill/bolundHill_water}
	\caption{Aerial photograph of the Bolund peninsula taken from a $125\,\mathrm{m}$ meteorological mast south of the site at Ris{\o} DTU. Figure adapted from \citet{berg2011bolund} under the Creative Commons Attribution-Noncommercial license.
	}
	\label{fig:realbolund}
\end{figure}
High-resolution topographic data obtained from airborne laser scanning and ground-based surveys reveal a steep western escarpment followed by a relatively flat plateau at the hilltop, as shown by the contour map in Fig.~\ref{fig:contourbolund}. The northern, southern, and eastern slopes are also steep, with local inclinations reaching up to $40^\circ$, although their crests are more rounded than the western face. Except for the escarpment and a narrow beach near the shoreline, the hill surface is uniformly covered with short grass, allowing it to be represented using a homogeneous aerodynamic roughness length in numerical models. Field observations report a roughness length of $z_0 = 0.015~\mathrm{m}$ for the hill surface, while the roughness length of the surrounding water is $z_0 = 0.0003~\mathrm{m}$ \cite{berg2011bolund}. Owing to its sharp geometric features, pronounced roughness transitions, and extensive experimental documentation, Bolund has become a canonical benchmark for the validation of atmospheric flow models over complex terrain.

\begin{figure}[htb!]

	\includegraphics[width=0.5\textwidth]{plots/bolundhill/bolundHill} 

	\caption{Terrain elevation map of the Bolund hill showing the contour lines of surface height. The red squares labelled M1--M8 indicate the locations of measurement probes (masts) distributed around the hill. The dashed lines show the two reference transects used in the Bolund experiment: Line B oriented at $270^\circ$ (westward) and Line A oriented at $239 ^\circ$ \cite{berg2011bolund}. Elevation is indicated by the color scale in metres.}

	\label{fig:contourbolund}
\end{figure}

The ABL flow over the Bolund terrain is simulated using the hybrid wall-modeled immersed boundary method to represent the complex hill geometry. The computational domain extends $L_x = 1000\,\mathrm{m}$, $L_y = 500\,\mathrm{m}$, and $L_z = 120\,\mathrm{m}$ in the streamwise, spanwise, and vertical directions, respectively, and is discretized using a uniform Cartesian grid of $400 \times 200 \times 112$ cells. Such grid resolution is relatively coarser than that employed in several previous studies, and reflects a practical LES resolution for ABL flow over complex terrain \cite{diebold2013flow,schubiger2020evaluation}.

Spatial derivatives of the convective terms are evaluated using a fourth-order central difference scheme augmented with a hyperviscous stabilization, while subgrid-scale stresses are modeled using the mixed Bardina-AMD (BAMD) formulation. The hybrid wall model is applied in conjunction with the immersed boundary framework to account for near-surface momentum transfer over the terrain. The aerodynamic roughness length of the hill surface is prescribed as $z_{0,\mathrm{ibm}} = 0.015\,\mathrm{m}$, consistent with a grass-covered surface, whereas the surrounding water surface is assigned a lower roughness length of $z_{0,\mathrm{water}} = 0.0003\,\mathrm{m}$ \cite{berg2011bolund}.

Inflow turbulent conditions are generated from a precursor simulation over flat terrain without the hill, driven by a constant streamwise pressure gradient. The pressure gradient is adjusted such that the resulting mean velocity profile follows the logarithmic law of the wall,
\begin{equation}
	U(z) = \frac{u_*}{\kappa} \ln\!\left(\frac{z}{z_{0,\mathrm{water}}}\right),
\end{equation}
where $u_*$ is the friction velocity, $\kappa$ is the von K\'arm\'an constant. The undisturbed reference velocity $U_\infty$ is defined as the mean velocity at the hill height obtained from the precursor simulation. Periodic boundary conditions are imposed in the spanwise direction, inflow-outflow conditions are applied in the streamwise direction, and a stress-free condition is enforced at the domain top.

The resolved three-component velocity field, $\boldsymbol{u}(t)$, is extracted from the LES using single-point probes placed at the mast locations along line~B (M7, M6, M3, and M8), as well as at additional probes positioned far upstream and downstream of the hill (see Fig.~\ref{fig:contourbolund}). At each probe, time series are recorded over the statistically stationary portion of the simulation with a sampling interval of $0.005\,\mathrm{s}$, corresponding to a sampling frequency of $200\,\mathrm{Hz}$. The simulation is advanced for a total of $20$ flow-through times to eliminate initial transients and establish a fully developed turbulent state. Statistics are then computed using the last $5$ flow-through times. The comparison between field measurements and LES is performed using the speed-up effect and turbulent kinetic energy (TKE) as key diagnostic quantities.

The terrain-induced speedup is quantified using the standard speed-up factor
\begin{equation}
	\Delta S = \frac{U - U_{\mathrm{ref}}}{U_{\mathrm{ref}}},
\end{equation}
where $U$ denotes the wind speed at a given mast location and height, and $U_{\mathrm{ref}}$ is the undisturbed upstream reference mean speed at the same height.
\begin{figure}[htbp]
	\begin{tabular}{c}
		(a) Speed up effect at 2m AGL \\
		\includegraphics[width=0.6\linewidth]{plots/bolundhill/speedupAGL2m} \\
		(b) Speed up effect at 5m AGL \\
		\includegraphics[width=0.6\linewidth]{plots/bolundhill/speedupAGL5m}\\
		
	\end{tabular}
	\caption{Comparison of the terrain-induced speedup factor $\Delta S$ along line~B for field measurements and LES predictions at (a)~2~m and (b)~5~m above ground level (a.g.l.). The speed-up is evaluated at the mast locations M7, M6, M3, and M8 and is defined relative to an upstream reference velocity. Field data are shown by blue circles, while LES results are shown by red squares. The horizontal dashed line indicates zero speedup. The shaded gray region represents a schematic projection of the Bolund hill profile along line~B for reference.
	}
	
	\label{fig:speedupBoulnd}
\end{figure}
Figure~\ref{fig:speedupBoulnd}(a-b) compares the speed-up factor $\Delta S$ obtained from LES and field measurements at (a)~2~m and (b)~5~m above ground level (a.g.l.). At 2~m a.g.l., the field measurements indicate negative speed-up at all mast locations (M7, M6, M3, and M8), reflecting an overall near-surface deceleration along line~B. In contrast, the LES reproduces the negative values at M7, M3, and M8 but predicts a positive speedup at M6, leading to the dominant mismatch at this windward-slope location. This discrepancy at M6 may be attributed to the strong sensitivity of the near-surface mean flow to the wall-modeled immersed boundary treatment and to local flow adjustment processes on the windward face, including the possible formation of a separation bubble, which are particularly challenging to capture \cite{schubiger2020evaluation}. At 5~m a.g.l., the agreement between LES and field measurements improves significantly at all mast locations. Both the magnitude and spatial variation of the speedup are captured with excellent fidelity, reflecting the reduced influence of the immediate surface layer and diminished sensitivity to near-wall modelling at this elevation. Overall, the results demonstrate that the LES reliably predicts the terrain-induced speedup over the Bolund hill, with remaining discrepancies largely confined to the near-surface windward-slope location at 2~m a.g.l.
\begin{figure}[htbp]
	\begin{tabular}{c}
		(a) Turbulent kinetic energy at 2m AGL \\
		\includegraphics[width=0.6\linewidth]{plots/bolundhill/tkeAGL2m} \\
		(b) Turbulent kinetic energy at 5m AGL \\
		\includegraphics[width=0.6\linewidth]{plots/bolundhill/tkeAGL5m}\\
		
	\end{tabular}
	\caption{Comparison of the turbulent kinetic energy (TKE) along line~B for field measurements and LES predictions at (a)~2~m and (b)~5~m above ground level (a.g.l.). The TKE is computed from the resolved velocity fluctuations and normalized by the square of the friction velocity, $u_*^2$. Values are shown at the mast locations M7, M6, M3, and M8. Field measurements are indicated by blue circles, and LES results are shown by red squares. The shaded gray region represents a schematic projection of the Bolund hill profile along line~B for reference.
	}
	
	\label{fig:tkeBoulnd}
\end{figure}
Figure~\ref{fig:tkeBoulnd}(a-b) compares the TKE obtained from LES and field measurements at (a)~2~m and (b)~5~m above ground level (a.g.l.) along line~B. The TKE is computed from the resolved velocity fluctuations as TKE~$=0.5\,(\langle {u'u'} \rangle +\langle {v'v'} \rangle +\langle{w'w'} \rangle )$, where $u'$, $v'$, and $w'$ denote the fluctuating components of the resolved velocity field. Only the resolved-scale contribution to the TKE is considered in this analysis; subgrid-scale contributions are omitted from the comparison.

At 2~m a.g.l., the field measurements exhibit a pronounced peak in TKE at the windward-slope location M6, while the remaining mast locations (M7, M3, and M8) show comparatively lower levels. The LES reproduces the measured TKE well at M7, M3, and M8, indicating good agreement away from the immediate windward-slope region. The primary discrepancy is confined to M6, where the LES underpredicts the magnitude of the TKE peak observed in the field data. This mismatch may be attributed to the strong sensitivity of near-surface turbulence production to the wall-modeled immersed boundary treatment, the difficulty in accurately capturing the separation bubble, and the omission of subgrid-scale contributions to the total turbulent kinetic energy. At 5~m a.g.l., the agreement between LES and field measurements improves considerably at all mast locations. Both datasets exhibit a smooth increase in TKE from upstream to downstream, and the LES reproduces both the magnitude and spatial variation of the measured TKE with good accuracy. The improved agreement at this height reflects the reduced influence of near-wall processes and subgrid-scale effects, indicating that the resolved turbulence is more reliably represented away from the immediate surface layer.

\begin{table}[htb!]
	\centering
	\renewcommand{\arraystretch}{1.2}
	\begin{tabular}{l
			c
			@{\hspace{1.5cm}}  % extra horizontal space
			c
			c}
		\hline
		\textbf{Model type} &
		\textbf{$\mathrm{MAE}_{\Delta S}[\%]$} &
		\textbf{$\mathrm{MAE}_{k}[\%]$} &
		\textbf{Reference} \\
		\hline
		RANS (2 eq.) & 15.1 & 47.0 & \citet{bechmann2011bolund} \\
		RANS (1 eq.) & 17.2 & 44.7 & \citet{bechmann2011bolund} \\
		Experiment   & 14.7 & 61.4 & \citet{bechmann2011bolund} \\
		LES          & 17.3 & 48.0 & \citet{bechmann2011bolund} \\
		Linearized model & 23.7 & 76.7 & \citet{bechmann2011bolund} \\
		RANS (2 eq.) & 10.35 & -- & \citet{prospathopoulos2012application} \\
		LES (EPFL)   & 12.10 & -- & \citet{diebold2013flow} \\
		LES          & 10.9 & 19.3 & \citet{vuorinen2015large} \\
		LBM–LES      & 8.0 & -- & \citet{schubiger2020evaluation} \\
		Fluent RANS  & 10.0 & -- &  \citet{schubiger2020evaluation} \\
		Fluent DES   & 17.3 & -- & \citet{schubiger2020evaluation}  \\
		Present  & 8.35 & 15.9 & - \\
		\hline
	\end{tabular}
	\caption{Comparison of mean absolute errors (MAE) for the terrain-induced speed-up factor, $\mathrm{MAE}_{\Delta S}$, and turbulent kinetic energy, $\mathrm{MAE}_{k}$, reported for different modeling approaches applied to the Bolund hill benchmark. The MAE values are expressed in percentage and are computed by comparing model predictions with field measurements at multiple mast locations and heights. Reported values are taken directly from the referenced studies where available; dashes indicate quantities that were not reported. The final row corresponds to the present hybrid wall-modeled immersed boundary and BAMD SGS model, with MAE values averaged over the 2~m and 5~m above ground level measurements.
	}
	\label{tab:meanabsoluteerror}
\end{table}
Next, the quantitative agreement between LES predictions and field measurements is assessed using the mean absolute error (MAE) computed for both the speedup factor and the turbulent kinetic energy. The MAE for the speed-up is defined as
\[
\mathrm{MAE}_{\Delta S}[\%]
= 100\,\frac{1}{N}\sum_{i=1}^{N}
\left|\Delta S_i^{\mathrm{LES}}-\Delta S_i^{\mathrm{field}}\right|,
\]
while the MAE for turbulent kinetic energy is evaluated as
\[
\mathrm{MAE}_{k}[\%]
= 100\,\frac{1}{N}\sum_{i=1}^{N}
\left|k_i^{\mathrm{LES}}-k_i^{\mathrm{field}}\right|,
\]
where \(i\) denotes the probe index along line~B and \(N\) is the number of valid probe locations included in the comparison. The resulting MAE values are summarized in Table~\ref{tab:meanabsoluteerror} together with corresponding values reported in the literature for a range of modeling approaches. The present study exhibits the second lowest average MAE for the speedup factor, with a value of \(8.35\%\) when averaged over the 2~m and 5~m a.g.l.\ levels, and also achieves the lowest reported MAE for turbulent kinetic energy at \(15.9\%\). It is important to note that a substantial fraction of the remaining error originates from near-surface discrepancies at 2~m a.g.l. where flow separation, strong shear, and wall-model limitations introduce additional uncertainty compared to measurements at higher elevations.

\section{\label{sec:conclusion}Conclusion}

A hybrid wall-modeled immersed boundary method combined with a mixed subgrid-scale closure was assessed for large eddy simulation of atmospheric boundary layer flow over complex terrain. The framework was evaluated using numerical simulations over flat terrain, an idealized three-dimensional hill, and a realistic terrain configuration corresponding to the Bolund hill benchmark. The simulations reproduce mean flow fields, turbulence intensities, and energy spectra with good agreement to experimental measurements and high-resolution numerical reference data, despite the use of relatively coarse grids. Terrain-induced effects, including flow deceleration, acceleration, separation, and recovery, are captured with consistent accuracy.

A comparative assessment of the tested subgrid-scale closures shows that most models exhibit a strong dependence on grid resolution, with noticeable variation in predicted turbulence statistics as the mesh is refined or coarsened. In contrast, the BAMD model remains comparatively consistent across resolutions. The results further indicate that the overall robustness of the predictions arises from the combined action of the wall-modeled immersed boundary method and the SGS closure rather than from either component alone. The wall-modeled immersed boundary treatment provides consistent near-surface stress over non grid-aligned terrain, while the mixed subgrid-scale closure preserves the fraction of total turbulent kinetic energy across resolutions. Their combined use reduces grid sensitivity and improves the reliability of second-order statistics in heterogeneous shear flows.

The current hybrid wall-modeled immersed boundary method is formulated only for neutral atmospheric boundary layer flows, so extending it to stable and unstable surface layers with appropriate stability corrections is a necessary next step. This will allow the method to be applied more reliably under realistic atmospheric stratification.

\begin{acknowledgements}
This research has been supported by the Natural Sciences and Engineering Research Council of Canada (grant no. 556326) in partnership with UL Renewables. We also gratefully acknowledged the computational resources provided by the Digital Research Alliance of Canada. 
\end{acknowledgements} 

% \nocite{*}

\bibliography{apssamp}% Produces the bibliography via BibTeX.

\end{document}
%
% ****** End of file apssamp.tex ******
